\exercisesection{2}

\begin{enumerate}
\def\labelenumi{\arabic{enumi}.}
\item
  Let K be a commutative field, A the polynomial ring K{[}X, Y{]} in two indeterminates and \(a_{n}\) the principal ideal \((XY^{n})\) in A; the sequence \((a_{n})_{n\geqslant1}\) forms with \(a_{0}=A\) an exhaustive and separated filtration on A. Let b be the principal ideal (X) of A; show that the topology on b induced by that on A is strictly coarser than the topology defined by the filtration on b associated with that on A (and a fortiori this latter filtration is distinct from the filtration induced by that on A).
\item
  Let K be a commutative field of characteristic \(\neq2\) and A the ring K{[}{[}X, Y{]}{]} of formal power series in two indeterminates.
\end{enumerate}

\begin{enumerate}
\def\labelenumi{(\alph{enumi})}
\item
  Show that in A the principal ideal \(\mathfrak{p} = (\mathrm{X}^2 -\mathrm{Y}^3)\) is prime. (If a product \(f(\mathbf{X},\mathbf{Y})g(\mathbf{X},\mathbf{Y})\) of two formal power series is divisible by \(\mathbf{X}^2 -\mathbf{Y}^3\) , note first that \(f(\mathrm{T}^3,\mathrm{T}^2) = 0\) or \(g(\mathrm{T}^3,\mathrm{T}^2) = 0\) in the ring of formal power series \(\mathbf{K}[[\mathrm{T}]]\) ; assuming for example \(f(\mathrm{T}^3,\mathrm{T}^2) = 0\) , show first that \(f(\mathbf{X},\mathbf{Y}^2)\) is divisible by \(\mathbf{X} - \mathbf{Y}^3\) and by \(\mathbf{X} + \mathbf{Y}^3\) and, by considering \(f(-\mathbf{Y}^3 +\mathbf{X},\mathbf{Y}^2)\) , prove finally that \(f(\mathbf{X},\mathbf{Y}^2)\) is divisible by \(\mathbf{X}^2 -\mathbf{Y}^6\) .)
\item
  Let \(m\) be the maximal ideal \(AX + AY\) of \(A\) . Show that \(p\) is closed with respect to the \(m\) -adic topology on \(A\) (with which \(A\) is Hausdorff and complete) but that in the ring \(\text{gr}(A)\) , \(\text{gr}(p)\) is not a prime ideal ( \(p\) being given the filtration induced by that on \(A\) ).
\end{enumerate}

\begin{enumerate}
\def\labelenumi{\arabic{enumi}.}
\setcounter{enumi}{2}
\item
  Let A be a filtered ring and E a finitely generated A-module. Show that, if E is given the filtration induced by that on A, gr(E) is a finitely generated gr(A)-module (cf.~Exercise 5 (c)).
\item
  Give an example of a bijective A-linear mapping \(u \colon \mathbf{E} \to \mathbf{F}\) , where \(\mathbf{E}\) and \(\mathbf{F}\) are two filtered A-modules, such that \(u\) is compatible with the filtrations but \(\operatorname{gr}(u)\) is neither injective nor surjective. (Take \(\mathbf{E} = \mathbf{A}\) , \(\mathbf{F} = \mathbf{A}\) with another filtration and \(u\) the identity mapping.)
\item
  Let A be a filtered commutative ring with filtration \((\mathfrak{a}_n)_{n\geqslant 0}\) such that \(\mathfrak{a}_0 = \mathbf{A}\) .
\end{enumerate}

\begin{enumerate}
\def\labelenumi{(\alph{enumi})}
\tightlist
\item
  For the ring gr(A) to be generated by a family of elements whose degrees are bounded, it is necessary and sufficient that there exist an integer q such that, for all \(n\) , \(\mathfrak{a}_n = \mathfrak{a}_{n+1} + \mathfrak{b}_n\) , where \(\mathfrak{b}_n\) is the sum of the ideals \(\mathfrak{a}_1^{\alpha_1}\mathfrak{a}_2^{\alpha_2}\ldots\mathfrak{a}_q^{\alpha_q}\) for
\end{enumerate}

\[
\alpha_ {i} \geqslant 0 (1 \leqslant i \leqslant q)
\]

\[
\sum_ {i = 1} ^ {q} i \alpha_ {i} = n.
\]

\[
\mathfrak {a} _ {n} = \mathfrak {a} _ {n + k} + \mathfrak {b} _ {n}
\]

\[
k > 0
\]

\begin{enumerate}
\def\labelenumi{(\alph{enumi})}
\setcounter{enumi}{1}
\item
  For \(\mathrm{gr}(\mathbf{A})\) to be Noetherian, it is necessary and sufficient that \(\mathbf{A} / \mathfrak{a}_1\) be Noetherian, that the condition of (a) hold and that for \(i \leqslant q\) the \(\mathfrak{a}_i / \mathfrak{a}_{i+1}\) be finitely generated \((\mathbf{A} / \mathfrak{a}_1)\) -modules (use Corollary 4 to Theorem 2 of no. 10).
\item
  Let K be a commutative field and A the polynomial ring in two indeterminates K{[}X, Y{]}. A total ordering is defined on the set of monomials \(X^{m}Y^{n}\) by setting \(X^{m}Y^{n} \leqslant X^{p}Y^{q}\) if \(m + n \leqslant p + q\) or if \(m + n = p + q\) and \(m \leqslant p\) . Let \((\mathbf{M}_{k})\) be the sequence of monomials in X, Y thus arranged in ascending order and let \(a_{n}\) be the ideal of A generated by the \(M_{k}\) of index \(k \geqslant n\) . Show that \((a_{n})\) is a filtration compatible with the ring structure on A; with this filtration gr(A) admits divisors of 0 and is not Noetherian, although A is a Noetherian integral domain (use the criterion in (b)); in particular, gr( \(a_{1}\) ) (with the filtration induced on \(a_{1}\) by that on A) is not a finitely generated gr(A)-module, although \(a_{1}\) is a finitely generated A-module (cf.~§ 1, no. 2, Corollary to Proposition 1).
\end{enumerate}

¶ 6. Let A be a commutative ring, E and F two A-modules and (E \(_{n}\) ) (resp. (F \(_{n}\) )) an exhaustive filtration on E (resp. F) consisting of sub-A-modules. On the tensor product G = E ⊗A F, consider the exhaustive filtration consisting of the G \(_{n}\) = ∑ \(_{i+j=n}\) Im(E \(_{i}\) ⊗A F \(_{j}\) ).

\begin{enumerate}
\def\labelenumi{(\alph{enumi})}
\tightlist
\item
  Show that the composite canonical homomorphisms
\end{enumerate}

\[
\begin{array}{r l}&(\mathrm{E} _ {i} / \mathrm{E} _ {i + 1}) \otimes_ {\mathrm{A}} (\mathrm{F} _ {j} / \mathrm{F} _ {j + 1}) \rightarrow\\&\qquad (\mathrm{E} _ {i} \otimes_ {\mathrm{A}} \mathrm{F} _ {j}) / (\operatorname{Im} (\mathrm{E} _ {i} \otimes_ {\mathrm{A}} \mathrm{F} _ {j + 1}) + \operatorname{Im} (\mathrm{E} _ {i + 1} \otimes_ {\mathrm{A}} \mathrm{F} _ {j})) \rightarrow \mathrm{G} _ {i + j} / \mathrm{G} _ {i + j + 1}\end{array}
\]

are the restrictions of a graded homomorphism of degree 0 (called canonical) \(\operatorname{gr}(\mathbf{E}) \otimes_{\mathbf{A}} \operatorname{gr}(\mathbf{F}) \to \operatorname{gr}(\mathbf{E} \otimes_{\mathbf{A}} \mathbf{F})\) , which is surjective.

\begin{enumerate}
\def\labelenumi{(\alph{enumi})}
\setcounter{enumi}{1}
\item
  Show that, if \(\operatorname{gr}(\mathbf{E})\) is a flat A-module, the A-modules \(\mathrm{E}_m / \mathrm{E}_n\) for \(m \leqslant n\) and \(\mathrm{E} / \mathrm{E}_n\) for \(n \in \mathbf{Z}\) are flat.
\item
  Suppose in the following that \(\mathrm{gr}(\mathbf{E})\) is a flat A-module and that \(\mathbf{E}\) is a flat A-module (the second hypothesis being a consequence of the first if \((\mathbf{E}_n)\) is a discrete filtration). Then the \(\mathrm{H}_{ij} = \mathrm{E}_i \otimes_{\mathbf{A}} \mathrm{F}_j\) are canonically identified with submodules of \(\mathbf{G} = \mathbf{E} \otimes_{\mathbf{A}} \mathbf{F}\) (Chapter I, § 2, no. 5, Proposition 4). Show that, for any two finite subsets \(\mathbf{R}\) , \(\mathbf{S}\) of \(\mathbf{Z} \times \mathbf{Z}\) ,
\end{enumerate}

\[
\left(\sum_ {(i, j) \in \mathbf {R}} \mathrm{H} _ {i j}\right) \cap \left(\sum_ {(h, k) \in \mathbf {S}} \mathrm{H} _ {h k}\right) = \mathrm{H} _ {\sup (i, h), \sup (j, k)}\tag{*}
\]

where in the second sum \((i,j,h,k)\) run through \(R \times S\) . (If R and S are each reduced to a single element, use Chapter I, §2, no. 6, Proposition 7. If p is the greatest of the indices i and h appearing in the elements of R or S, q the least of the indices j and k appearing in the elements of R or S, consider the images of the two sides of (*) under the canonical homomorphism

\[
\mathbf {E} \otimes_ {\mathbf {A}} \mathbf {F} _ {q} \rightarrow (\mathrm{E} / \mathrm{E} _ {p}) \otimes_ {\mathbf {A}} \mathbf {F} _ {q}
\]

and argue by induction on \(\operatorname{Card}(\mathbf{R}) + \operatorname{Card}(\mathbf{S})\) .

\begin{enumerate}
\def\labelenumi{(\alph{enumi})}
\setcounter{enumi}{3}
\tightlist
\item
  Deduce from (c) that the canonical homomorphism defined in (a) is then bijective.
\end{enumerate}

\begin{enumerate}
\def\labelenumi{\arabic{enumi}.}
\setcounter{enumi}{6}
\tightlist
\item
  Let A be a commutative ring, E an A-module, F a submodule of E, B the exterior algebra \(\wedge\) (E) and \(\Im\) the two-sided ideal of B generated by the canonical images in B of the elements of F. Let \(\operatorname{gr}^{\Im}(B)\) be the graded ring associated with the ring B filtered by the \(\Im\) -adic filtration.
\end{enumerate}

\begin{enumerate}
\def\labelenumi{(\alph{enumi})}
\item
   Define a surjective canonical (non-graded) A-algebra homomorphism \((\wedge (\mathbf{F})) \otimes_{\mathbf{A}}^{g}(\wedge (\mathbf{E} / \mathbf{F})) \to \mathrm{gr}^{\Im} (\mathbf{B})\) (where we mean the skew tensor product of graded algebras (Algebra, Chapter III).
\item
  Show that, if \(\mathbf{F}\) admits a complement in \(\mathbf{E}\) , the homomorphism defined in (a) is an isomorphism.
\item
   Take \(\mathbf{A} = \mathbf{Z}\) , \(\mathbf{E} = \mathbf{Z} / 4\mathbf{Z}\) , \(\mathbf{F} = 2\mathbf{Z} / 4\mathbf{Z}\) ; show that the homomorphism defined in (a) is not injective in this case.
\end{enumerate}

\begin{enumerate}
\def\labelenumi{\arabic{enumi}.}
\setcounter{enumi}{7}
\item
  Let A be a filtered ring, \((\mathbf{A}_{n})\) its filtration, E a filtered A-module and \((\mathbf{E}_{n})\) its filtration. Show that if \(A_{0} = A\) and \(E_{0} = E\) , the mapping \((a, x) \mapsto ax\) of \(A \times E\) to E is uniformly continuous with the topologies defined by the filtrations.
\item
  Give an example of two filtered rings A, B whose filtrations are exhaustive and separated and a non-surjective homomorphism \(u: A \to B\) compatible with the filtrations and such that \(\text{gr}(u)\) is bijective. From this deduce a counterexample to Corollary 1 to Proposition 12 when the ring A is not assumed to be complete and a counter-example to Proposition 13 when E is no longer assumed to be finitely generated (use Algebra, Chapter VIII, §7, Exercise 3(b)).
\item
  Let A be a Noetherian ring and \(\sigma\) an automorphism of A. Show that the ring E defined in Algebra, Chapter IV, § 5, Exercise 10 (b) is a (left or right) Noetherian ring.
\item
  Show that, if E is a vector space of dimension \(\geqslant2\) over a commutative field, the tensor algebra of E is a ring which is neither left nor right Noetherian (if a, b are two linearly independent vectors in E, consider the left ideal (or right ideal) generated by the elements \(a^{n}b^{n}\) for \(n\geqslant1\) ).
\end{enumerate}

 ¶ 12. Let K be a commutative field, A the polynomial ring \(K[X_{i}]_{i\in I}\) in an arbitrary infinite family of indeterminates and m the (maximal) ideal of A generated by the \(X_{i}\) . If we set \(A_{i}=A/m^{i+1}\) , the \(A_{i}\) and the canonical homomorphisms \(h_{ij}:A/m^{j+1}\to A/m^{i+1}\) for \(i\leqslant j\) satisfy the conditions of Proposition

14; the ring \(\lim A_{i}\) is the completion \(\hat{A}\) of \(A\) with respect to the m-adic topology and the kernel of the canonical homomorphism \(\hat{A} \to A_{i}\) is equal to \((m^{i})^{\wedge}\) , the closure of \(m^{i}\) in \(\hat{A}\) .

\begin{enumerate}
\def\labelenumi{(\alph{enumi})}
\item
  Show that \(\hat{\mathbf{A}}\) is canonically identified with the ring of formal power series in the \(\mathbf{X}_{\iota}\) , each of which has only a finite number of terms of given degree (Algebra, Chapter IV, § 5, Exercise 1).
\item
  From now on take \(\mathbf{I} = \mathbf{N}\) . Show that \(\hat{\mathfrak{m}} \neq \hat{\mathbf{A}}. \mathfrak{m}\) (consider the formal power series \(\sum_{n=1}^{\infty} \mathbf{X}_n^n\) ).
\item
  Suppose that K is a finite field. Show that \((\hat{\mathfrak{m}})^{2} \neq (\mathfrak{m}^{2})^{\bullet}\) . (Show first the following result: for every integer k \textgreater{} 0, there exists an integer \(n_{k}\) and for every integer \(n \geqslant n_{k}\) a homogeneous polynomial \(F_{n}\) of degree n in \(n^{2}\) indeterminates with coefficients in K such that \(F_{n}\) cannot be the sum of the terms of degree n in any polynomial of the form \(P_{1}Q_{1} + \cdots + P_{k}Q_{k}\) , where the \(P_{i}\) and \(Q_{i}\) are polynomials without constant term in the same \(n^{2}\) indeterminates).
\end{enumerate}

Deduce that \(\hat{A}\) is not complete with respect to the \(\hat{m}\) -adic topology.

\begin{enumerate}
\def\labelenumi{\arabic{enumi}.}
\setcounter{enumi}{12}
\tightlist
\item
  Let K be a field, \(A = K[[X]]\) the ring of formal power series and m its maximal ideal, so that A is Hausdorff and complete with the m-adic topology (no. 6, Corollary to Proposition 6). On the additive group A consider the filtration \((\mathrm{E}_{n})\) such that \(E_{0} = A\) and \(E_{n}\) is the intersection of \(m^{n}\) and the ring K{[}X{]}; this filtration is exhaustive and separated, the topology T which it defines on A is compatible with the additive group structure on A and is finer than the m-adic topology but A is not a complete group with the topology T (consider the sequence of polynomials \((1 - X^{n})/(1 - X)\) ).
\end{enumerate}

14 Let A be a ring (not necessarily commutative) and E a left A-module. A topology on E is called linear if it is invariant under translations and 0 admits a fundamental system of neighbourhoods which are submodules of E; E is then said to be linearly topologized. A linear topology on E is compatible with its additive group structure and defines on E a topological A-module structure if A is given the discrete topology. On any A-module the discrete topology and the coarsest topology are linear topologies.

\begin{enumerate}
\def\labelenumi{(\alph{enumi})}
\item
  If E is a linearly topologized module and F a submodule of E, the induce topology on F and the quotient topology of that on E by F are linear topologies. If \((\mathrm{E}_{\alpha}, f_{\alpha\beta})\) is an inverse system of linearly topologized A-modules, where the \(f_{\alpha\beta}\) are continuous linear mappings, the topological A-module \(E = \lim_{\leftarrow} E_{\alpha}\) is linearly topologized.
\item
  Let E be a linearly topologized A-module. There exists a fundamental system \((\mathrm{V}_{\lambda})_{\lambda \in L}\) of neighbourhoods of 0 in E consisting of open (and closed) submodules; if \(\phi_{\lambda\mu}: E/V_{\mu} \to E/V_{\lambda}\) is the canonical mapping when \(V_{\lambda} \supset V_{\mu}\) , the family \((\mathrm{E}/\mathrm{V}_{\lambda}, \phi_{\lambda\mu})\) is an inverse system of discrete A-modules; the topological A-module \(\tilde{E} = \lim_{\leftarrow} E/V_{\lambda}\) is identified with the Hausdorff completion of E, which is therefore linearly topologized (General Topology, Chapter III, § 7, no. 3).
\item
  A Hausdorff linearly topologized A-module E is discrete if it is Artinian or if there exists a least element in the set of submodules \(\neq0\) of E.
\end{enumerate}

\begin{enumerate}
\def\labelenumi{\arabic{enumi}.}
\setcounter{enumi}{14}
\tightlist
\item
  \begin{enumerate}
  \def\labelenumii{(\alph{enumii})}
  \tightlist
  \item
    Let E be a linearly topologized A-module (Exercise 14). For a filter base B on E consisting of linear varieties to admit a cluster point, it is necessary and sufficient that there exists a convergent filter base \(B' \supset B\) consisting of linear affine varieties. E is said to be linearly compact if it is Hausdorff and every filter base on E consisting of affine linear varieties admits at least one cluster point. Every Artinian module is linearly compact with the discrete topology. Every linearly compact submodule of a Hausdorff linearly topologized module F is closed in F.
  \end{enumerate}
\end{enumerate}

\begin{enumerate}
\def\labelenumi{(\alph{enumi})}
\setcounter{enumi}{1}
\item
  If E is a linearly compact A-module and u a continuous linear mapping of E to a Hausdorff linearly topologized A-module F, \(u(\mathbf{F})\) is a linearly compact submodule of F.
\item
  Let E be a Hausdorff linearly topologized A-module and F a closed submodule of E. For E to be linearly compact, it is necessary and sufficient that F and E/F be so.
\item
  Every product of linearly compact modules is linearly compact (consider a filter base which is maximal among the filter bases consisting of affine linear varieties.
\item
  Let \((\mathrm{E}_{\alpha}, f_{\alpha \beta})\) be an inverse system of linearly topologized modules relative to a directed indexing set; suppose that the \(f_{\alpha \beta}\) are continuous linear mappings and that, for \(\alpha \leqslant \beta\) , \(f_{\alpha \beta}^{-1}(0)\) is a linearly compact submodule of \(\mathbf{E}_{\beta}\) . Let \(\mathbf{E} = \varprojlim \mathbf{E}_{\alpha}\) and let \(f_{\alpha}\) be the canonical mapping \(\mathbf{E} \to \mathbf{E}_{\alpha}\) ; show that, for all \(\alpha, f_{\alpha}(\mathbf{E}) = \bigcap_{\alpha \leqslant \beta} f_{\alpha \beta}(\mathbf{E}_{\beta})\) (in particular, if the \(f_{\alpha \beta}\) are surjective, so are the \(f_{\alpha}\) ) and \(f_{\alpha}^{-1}(0)\) is linearly compact (use (d) and General Topology, Chapter I, Appendix, no. 2, Theorem 1).
\end{enumerate}

¶ 16. (a) Let E be a Hausdorff linearly topologized A-module. Show that the following properties are equivalent:

(α) E is linearly compact (Exercise 15).

(β) For every continuous linear mapping \(u\) from \(E\) to a Hausdorff linearly topologized A-module \(F\) , \(u(E)\) is a closed submodule of \(F\) .

\begin{enumerate}
\def\labelenumi{(\Alph{enumi})}
\setcounter{enumi}{24}
\tightlist
\item[(γ)]
  With every linear topology (Hausdorff or not) on E coarser than the given topology, E is complete.
\end{enumerate}

(δ) E is complete and there is a fundamental system \((\mathbf{U}_{\lambda})\) of open neighbourhoods of 0 in E, consisting of submodules and such that the \(E/U_{\lambda}\) are linearly compact discrete A-modules (cf.~§ 3, Exercise 5).

(To see that \((\gamma)\) implies \((\delta)\) consider an open submodule \(F\) of \(E\) and a filter base \(\mathfrak{B}\) on \(E / F\) consisting of affine linear varieties. For all \(V \in \mathfrak{B}\) , let \(M_V\) be the inverse image in E of the direction submodule of V in E/F; consider on E the linear topology in which the \(M_{V}\) form a fundamental system of neighbourhoods of 0.)

\begin{enumerate}
\def\labelenumi{(\alph{enumi})}
\setcounter{enumi}{1}
\item
  Let E be a Hausdorff linearly topologized A-module and M a linearly compact submodule of E. Show that, for every closed submodule F of E, \(M + F\) is closed in E (consider the image of M in E/F).
\item
  Let E be a linearly compact A-module and F a Hausdorff linearly topologized A-module. Show that for ever closed submodule M of \(E \times F\) , the projection of M onto F is closed. Obtain the converse (cf.~General Topology, Chapter I, § 10, no. 2).
\item
  Let E be a linearly compact A-module and u a continuous linear mapping of E to a Hausdorff linearly topologized A-module F. Show that, for every filter base B on E consisting of affine linear varieties, the image under u of the set of cluster points of B is the set of cluster points of \(u(\mathfrak{B})\) . In particular, for every closed submodule M of E, \(\bigcap_{N\in\mathfrak{B}}(M+N)=M+\bigcap_{N\in\mathfrak{B}}N\) .
\end{enumerate}

\begin{enumerate}
\def\labelenumi{\arabic{enumi}.}
\setcounter{enumi}{16}
\tightlist
\item
  A linear topology \(\mathcal{T}\) on an A-module E is called minimal if it is Hausdorff and there exists no Hausdorff linear topology strictly coarser than \(\mathcal{T}\) .
\end{enumerate}

\begin{enumerate}
\def\labelenumi{(\alph{enumi})}
\item
  For a Hausdorff linear topology T on E to be minimal, it is necessary and sufficient that every filter base B on E, consisting of affine linear varieties and having a single cluster point, be convergent to this point. (To see that the condition is necessary, observe that, when M runs through B and V a fundamental system of neighbourhoods of 0 consisting of submodules, the M + V form a filter base with the same cluster points as B. To see that the condition is sufficient, note that a filter base consisting of open submodules, whose intersection is reduced to 0, is a fundamental system of neighbourhoods of 0 in a Hausdorff linear topology coarser than T.)
\item
  For a discrete topology on an A-module E to be minimal, it is necessary and sufficient that E be Artinian.
\item
  If \(\mathcal{T}\) is minimal, the topology induced by \(\mathcal{T}\) on any closed submodule of \(\mathbf{E}\) is minimal.
\end{enumerate}

¶ 18. In an A-module E, a submodule M ≠ E is called sheltered if there exists a least element in the set of submodules ≠0 in E/M.

\begin{enumerate}
\def\labelenumi{(\alph{enumi})}
\item
  Show that every submodule \(N \neq E\) of E is an intersection of sheltered submodules (for all \(x \notin N\) , consider a maximal element in the set of submodules of E containing N and not containing x).
\item
  Let \(\mathcal{T}\) be a Hausdorff linear topology on E and let \(\mathcal{T}^*\) be the linear topology with fundamental system of neighbourhoods of 0 the filter base generated by the submodules of E which are open under \(\mathcal{T}\) and sheltered. Show that \(\mathcal{T}^*\) is Hausdorff and that every submodule which is closed under \(\mathcal{T}\) is also closed under \(\mathcal{T}^*\) (note that every submodule which is closed under \(\mathcal{T}\) is an intersection of submodules which are open under \(\mathcal{T}\) ). Deduce that, if E is complete with respect to \(\mathcal{T}^*\) , it is complete with respect to \(\mathcal{T}\) (General Topology, Chapter III, § 3, no. 5, Proposition 9).
\item
  Suppose that \(\mathcal{T}\) is linearly compact. Then show that \(\mathcal{T}^*\) is linearly compact and is the coarsest Hausdorff linear topology coarser than \(\mathcal{T}\) ; in particular \(\mathcal{T}^*\) is a minimal linear topology (Exercise 17). (Let \(\mathfrak{B}\) be a filter base consisting of submodules closed under \(\mathcal{T}\) whose intersection is 0; show that, for every submodule U which is open under \(\mathcal{T}\) and sheltered, there exists \(\mathbf{M} \in \mathfrak{B}\) such that \(\mathbf{M} \subset \mathbf{U}\) , using Exercise 16 (d)).
\item
  Let \(\mathbf{F}\) be a Hausdorff linearly topologized A-module and \(\mathcal{T}_1\) its topology; show that, if \(u: \mathrm{E} \to \mathrm{F}\) is a continuous linear mapping with respect to the topologies \(\mathcal{T}\) and \(\mathcal{T}_1\) , \(u\) is also continuous with respect to the topologies \(\mathcal{T}^*\) and \(\mathcal{T}_1^*\) .
\end{enumerate}

¶ 19. (a) Let E be a linearly compact A-module. Show that the following conditions are equivalent:

(α) For every continuous linear mapping u from E to a Hausdorff linearly topologized A-module F, u is a strict morphism (General Topology, Chapter III, § 2, no. 8) from E to F.

(β) For every closed submodule \(\mathbf{F}\) of \(\mathbf{E}\) , the quotient topology on \(\mathbf{E} / \mathbf{F}\) is a minimal topology (Exercise 17).

(γ) E is complete and there is a fundamental system \((\mathbf{U}_{\lambda})\) of open neighbourhoods of 0 in E, consisting of submodules and such that the \(E/U_{\lambda}\) are Artinian modules.

(To see that \((\gamma)\) implies \((\beta)\) , reduce it to the case where \(F = 0\) and use Exercises 17 and 18.)

If E satisfies the equivalent conditions \((\alpha)\) , \((\beta)\) and \((\gamma)\) , E is called strictly linearly compact.

\begin{enumerate}
\def\labelenumi{(\alph{enumi})}
\setcounter{enumi}{1}
\item
  Let E be a Hausdorff linearly topologized A-module and F a closed submodule of E. For E to be strictly linearly compact, it is necessary and sufficient that F and E/F be so.
\item
  Every inverse limit of strictly linearly compact modules is strictly linearly compact.
\item
  Let E be a Hausdorff linearly topologized A-module and u a linear mapping from E to a strictly linearly compact A-module F. Show that if the graph of u is closed in \(E \times F\) , u is continuous (use Exercise 16 (c) and the fact that, if M is closed in E and E/M is Artinian, M is open in E).
\end{enumerate}

¶ 20. (a) Show that a discrete linearly compact module cannot be the direct sum of an infinity of submodules which are not reduced to 0.

\begin{enumerate}
\def\labelenumi{(\alph{enumi})}
\setcounter{enumi}{1}
\item
  Give an example of a minimal linear topology (Exercise 17) which is not linearly compact (consider an infinite direct sum of simple modules no two of which are isomorphic).
\item
  Let E be a linearly compact module such that there exists a family of simple submodules whose sum is dense in E. Show that: (1) E is strictly linearly compact (apply (a) to the quotient of E by an open submodule); (2) E is isomorphic to the (topological) product of a family of discrete simple modules. (Consider the sets O of maximal open submodules of E such that for every finite
\end{enumerate}

\[
\mathrm{E} / \left(\bigcap_ {k = 1} ^ {n} \mathrm{G} _ {k}\right)
\]

sequence \((\mathbf{G}_k)_{1\leqslant k\leqslant n}\) of \(n\) distinct elements of \(\mathfrak{O}\) , \(\operatorname{E} / \left(\bigcap_{k=1}^{\infty}\mathbf{G}_k\right)\) is the direct sum of \(n\) simple submodules. Show that there exists a maximal set \(\mathfrak{O}\) (with respect to inclusion) and that the intersection of all the submodules \(\mathbf{G}\) belonging to such a set is reduced to 0; use the fact that \(\{0\}\) is an intersection of maximal open submodules and that for every maximal open submodules \(\mathbf{G}\) of \(\mathbf{E}\) there is at least one simple submodule of \(\mathbf{E}\) not contained in \(\mathbf{G}\) . Conclude by using the fact that \(\mathbf{E}\) is strictly linearly compact.)

\begin{enumerate}
\def\labelenumi{(\alph{enumi})}
\setcounter{enumi}{3}
\tightlist
\item
  Deduce from (c) that every linearly compact vector space over a field \(\mathbf{K}\) is strictly linearly compact and isomorphic to a product \(\mathbf{K}_s^{\mathrm{I}}\) .
\end{enumerate}

¶ 21. Recall that on a ring A a topology is called (left) linear if it is a linear topology on the A-module \(\mathbf{A}_s\) and it is compatible with the ring structure on A (Chapter II, § 2, Exercise 16). A topological ring A is called (left) linearly compact (resp. (left) strictly linearly compact) if \(\mathbf{A}_s\) is a linearly compact (resp. strictly linearly compact) A-module.

\begin{enumerate}
\def\labelenumi{(\alph{enumi})}
\item
  Show that, if A is left linearly compact with the topology T, the topology \(T^{*}\) defined in Exercise 18 (b) is also compatible with the ring structure on A (use Exercise 18 (d)).
\item
  Suppose that \(\mathbf{A}\) is commutative; let \(u \neq 0\) be an idempotent of \(\mathbf{A} / \Re\) , where \(\Re\) is the Jacobson radical of \(\mathbf{A}\) . Show that if \(\mathbf{A}\) is linearly compact with the discrete topology, there exists a unique idempotent \(e \in \mathbf{A}\) whose image in \(\mathbf{A} / \Re\) is equal to \(u\) . (Show that among the linear affine varieties \(x + \mathfrak{b}\) , where \(\mathfrak{b} \subset \Re\) and \(x^2 - x \in \mathfrak{b}\) , which are contained in the class \(u\) , there is a minimal element \(e + \mathfrak{a}\) . Deduce that \(e^2 = e\) and \(\mathfrak{a} = 0\) , by considering the element \(e^2 - e = r\) and showing that \(r \in \mathrm{Ar}^2\) , using Algebra, Chapter VIII, § 6, Exercise 10 (a). Prove the uniqueness of \(e\) by showing that if \(e_1, e_2\) are two idempotents of \(\mathbf{A}\) such that \(e_1 e_2 \in \Re\) , then \(e_1 e_2 = 0\) .)
\item
  Show that every commutative ring A which is linearly compact with the discrete topology is a direct composition of a finite number of local rings (which are linearly compact with the discrete topology). (Consider first the case, where \(\Re = 0\) , using Exercise 15 (b) and Chapter II, § 1, no. 2, Proposition 5, then apply (b) to the idempotents of A/ \(\Re\) .)
\item
  Show that every linearly compact (resp. strictly linearly compact) commutative ring A is the product of a family of linearly compact (resp. strictly linearly compact) local rings. (Let \((\mathfrak{m}_{\lambda})\) be the family of open maximal ideals of A; for every open ideal \(a\) not containing \(\mathfrak{m}_{\lambda}\) , let \(\mathrm{A}_{\lambda}^{\prime} \times \mathrm{A}_{\lambda}^{\prime \prime}\) be the decomposition of A/a as a direct composition of two rings such that \(\mathrm{A}_{\lambda}^{\prime}\) is the local ring with maximal ideal \((\mathfrak{m}_{\lambda} + \mathfrak{a}) / a\) defined in (c); let \(e_{\lambda}^{\prime}(a)\) be the unit element of \(\mathrm{A}_{\lambda}^{\prime}\) considered as a class mod. \(a\) in A; the \(e_{\lambda}^{\prime}(a)\) form a filter base which converges in A to an idempotent \(e_{\lambda}\) and \(Ae_{\lambda}\) is a closed ideal of A of which \(e_{\lambda}\) is the unit element. Show that the ring \(Ae_{\lambda}\) is a local ring and that A is isomorphic to the product of the \(Ae_{\lambda}\) .
\end{enumerate}

¶ 22. (a) Let A be a local ring and m its maximal ideal. Show that, if, for a linear topology T on A, A is strictly linearly compact, T is coarser than the m-adic topology.

\begin{enumerate}
\def\labelenumi{(\alph{enumi})}
\setcounter{enumi}{1}
\item
  Let A be a commutative ring and m an ideal of A. For A to be strictly linearly compact with the m-adic topology, it is necessary and sufficient that it be Hausdorff and complete with this topology, that A/m be an Artinian ring and \(m/m^{2}\) and (A/m)-module of finite length * (in other words, A is a complete Noetherian semi-local ring) *.
\item
  Let I be an infinite indexing set, K a finite field and \(A = K[[X_{t}]]_{t \in I}\) the algebra of formal power series in the family of indeterminates \((\mathbf{X}_{t})_{t \in I}\) (Algebra, Chapter IV, § 5, Exercise 1); A is a local ring. As a vector space over K, A may be identified with the product space \(\mathbf{K}^{\mathbf{N}^{(1)}}\) ; if K is given the discrete topology and A the product topology T, show that T is a linear topology on the ring A with which A is strictly linearly compact; if m is the maximal ideal of A, show, by arguing as in Exercise 12 (c), that A is not complete with the m-adic topology.
\end{enumerate}

¶ 23. Let A be a commutative ring with a linear topology T.

\begin{enumerate}
\def\labelenumi{(\alph{enumi})}
\item
  Show that the ideals \(\mathfrak{a}\) of \(\mathbf{A}\) , which are closed under \(\mathcal{T}\) and such that \(\mathbf{A} / \mathfrak{a}\) is Artinian, form a fundamental system of neighbourhoods of 0 under a linear topology \(\mathcal{T}^{(c)}\) on \(\mathbf{A}\) , which is coarser than \(\mathcal{T}\) (use Algebra, Chapter VIII, §2, Exercise 12); then \((\mathcal{T}^{(c)})^{(c)} = \mathcal{T}^{(c)}\) . Let \(\mathbf{B}\) be the Hausdorff completion ring of \(\mathbf{A}\) with respect to the topology \(\mathcal{T}^{(c)}\) ; show that \(\mathbf{B}\) is strictly linearly compact; \(\mathbf{B}\) is said to be the strictly linearly compact ring associated with \(\mathbf{A}\) .
\item
  Let \(\mathcal{T}_{\omega}(\mathbf{A})\) be the discrete topology on \(\mathbf{A}\) , \(\mathcal{T}_c(\mathbf{A})\) the topology \((\mathcal{T}_{\omega}(\mathbf{A}))^{(c)}\) . For the topology \(\mathcal{T}^{(c)}\) on \(\mathbf{A}\) to be Hausdorff, it is necessary and sufficient that \(\mathcal{T}\) be Hausdorff and that, for every ideal \(\mathfrak{b}\) of \(\mathbf{A}\) open under \(\mathcal{T}\) , the topology \(\mathcal{T}_c(\mathbf{A}/\mathfrak{b})\) be Hausdorff. * If \(\mathbf{A}\) is the ring of a non-discrete valuation of height 1 (Chapter VI), \(\mathcal{T}_c(\mathbf{A})\) is not Hausdorff. *
\item
  Suppose that A is linearly compact under \(\mathcal{T}\) ; then A is complete (but in general not Hausdorff) with \(\mathcal{T}^{(c)}\) . Then for \(\mathcal{T}^{(c)}\) to be Hausdorff, it is necessary and sufficient that there exist on A a linear topology which is strictly linearly compact and coarser than \(\mathcal{T}\) ; \(\mathcal{T}^{(c)}\) is then the unique topology with these properties.
\item
  Let E be a linearly topologized and strictly linearly compact A-module. Show that, if A is given the topology \(\mathcal{T}_{c}(A)\) , E is a topological ring (note that if F is an Artinian A-module, then, for all \(x \in F\) , the annihilator a of x is such that A/a is Artinian). Deduce that, if B is the Hausdorff completion of A with respect to the topology \(\mathcal{T}_{c}(A)\) , E can be considered as a topological B-module; the ring B being isomorphic to a product of strictly linearly compact local rings
\end{enumerate}

\(B_{\lambda}\) (Exercise 21 (d)), show that E is isomorphic to a product of strictly linearly compact submodules \(E_{\lambda}\) , where \(E_{\lambda}\) is annihilated by the \(B_{\mu}\) of index \(\mu \neq \lambda\) and can therefore be considered as a topological \(B_{\lambda}\) -module (if \(e_{\lambda}\) is the unit element of B, take \(E = e_{\lambda}.E\) ).

\begin{enumerate}
\def\labelenumi{\arabic{enumi}.}
\setcounter{enumi}{23}
\tightlist
\item
  On a commutative ring A let \(\mathcal{T}_{m}(A)\) denote the linear topology which has a fundamental system of neighbourhoods of 0 the products (equal to the intersections) of a finite number of powers of maximal ideals (cf.~no. 13, Proposition 17); if every finite intersection of ideals \(\neq0\) of A is \(\neq0\) , let \(\mathcal{T}_{u}(A)\) denote the linear topology for which a fundamental system of neighbourhoods of 0 consists of all the ideals \(\neq\{0\}\) of A.
\end{enumerate}

\begin{enumerate}
\def\labelenumi{(\alph{enumi})}
\item
  Show that \(\mathcal{T}_c(\mathbf{A})\) (Exercise 23) is coarser than \(\mathcal{T}_m(\mathbf{A})\) (observe that the Jacobson radical of an Artinian ring is nilpotent); give an example where \(\mathcal{T}_c(\mathbf{A}) \neq \mathcal{T}_m(\mathbf{A})\) (cf.~Exercise 22 (c)). * If A is the ring of a non-discrete valuation of height 1 (Chapter VI), then \(\mathcal{T}_c(\mathbf{A}) = \mathcal{T}_m(\mathbf{A})\) and \(\mathcal{T}_u(\mathbf{A}) \neq \mathcal{T}_m(\mathbf{A})\) . *
\item
  The topology \(\mathcal{T}_m(A)\) is the coarsest linear topology on \(A\) under which the maximal ideals of \(A\) are closed and every power of an open ideal is an open ideal (note that under a linear topology on \(A\) , under which a maximal ideal \(m\) is closed, \(m\) is necessarily open).
\item
  If A is a Noetherian ring, then \(\mathcal{T}_{m}(\mathbf{A}) = \mathcal{T}_{c}(\mathbf{A})\) (observe that, if \(m\) is a maximal ideal of A, \(\mathbf{A} / m^{k}\) is an Artinian ring for every integer \(k \geqslant 1\) , noting that \(\mathfrak{m}^h / \mathfrak{m}^{h+1}\) is an (A/m)-module which is necessarily of finite length). Give an example of a Noetherian local ring A for which \(\mathcal{T}_{m}(\mathbf{A}) \neq \mathcal{T}_{u}(\mathbf{A})\) (consider a local ring of a polynomial ring over a field).
\end{enumerate}

\begin{enumerate}
\def\labelenumi{\arabic{enumi}.}
\setcounter{enumi}{24}
\item
  Let A be a topological ring and E a finitely generated left A-module. Show that there exists on E a topology (called canonical) compatible with its A-module structure and finer than all the others (write E in the form \(A_{s}^{n}/R\) , give \(A_{s}^{n}\) the product topology and E the quotient topology). If \(E'\) is a topological A-module and \(u: E \to E'\) an A-linear mapping, show that u is continuous with the canonical topology on E. If further \(E'\) is finitely generated and has the canonical topology, u is a strict morphism. If the topology on A is defined by a filtration ( \(A_{n}\) ), show that the canonical topology on E is defined by the filtration ( \(A_{n}E\) ).
\item
  Let A be a commutative ring, m an ideal of A, \((a_{\lambda})_{\lambda \in L}\) a system of generators of m and \(\hat{A}\) the Hausdorff completion of A with respect to the m-adic topology.
\end{enumerate}

Show that, if \(\mathbf{A}'\) denotes the ring of formal power series in a family \((\mathrm{T}_{\lambda})_{\lambda \in \mathbf{L}}\) of indeterminates with only a finite number of terms of given degree (Exercise 12), \(\hat{\mathbf{A}}\) is isomorphic to the quotient of \(\mathbf{A}'\) by the closure \(\mathfrak{b}\) of the ideal of \(\mathbf{A}'\) generated by the \(T_{\lambda} - a_{\lambda}\) (use Theorem 1 of no. 8). In particular, the ring \(Z_{p}\) of p-adic integers is isomorphic to the quotient \(\mathbf{Z}[[T]]/(T-p)\) .

¶ 27. (a) Let A be a commutative ring with a linear topology. For every multiplicative subset S of A, let \(\mathrm{A}\{\mathrm{S}^{-1}\}\) denote the Hausdorff completion of the ring \(\mathrm{S}^{-1}\mathrm{A}\) with the topology for which a fundamental system of neighbourhoods of 0 consists of the ideals \(\mathrm{S}^{-1}\mathrm{U}_{\lambda}\) , where \((\mathrm{U}_{\lambda})\) is a fundamental system of neighbourhoods of 0 in A consisting of ideals of A. Show that \(\mathrm{A}\{\mathrm{S}^{-1}\}\) is canonically isomorphic to the inverse limit of the rings \(\mathrm{S}_{\lambda}^{-1}(\mathrm{A}/\mathrm{U}_{\lambda})\) where \(\mathrm{S}_{\lambda}\) is the canonical image of S in \(\mathrm{A}/\mathrm{U}_{\lambda}\) . If \(\hat{\mathrm{A}}\) is the Hausdorff completion of A, \(\mathrm{A}\{\mathrm{S}^{-1}\}\) is canonically isomorphic to \(\hat{\mathrm{A}}\{\mathrm{S}'^{-1}\}\) , where \(\mathrm{S}'\) is the canonical image of S in \(\hat{\mathrm{A}}\) .

\begin{enumerate}
\def\labelenumi{(\alph{enumi})}
\setcounter{enumi}{1}
\item
  For \(A\{S^{-1}\}\) to be reduced to 0, it is necessary and sufficient that, in A, 0 belong to the closure of S. Deduce an example where A is Hausdorff and complete but \(S^{-1}A\) is not so with the topology defined in (a).
\item
  For every continuous homomorphism \(u\) of \(A\) to a complete Hausdorff linearly topologized ring \(B\) such that \(u(S)\) consists of invertible elements of \(B\) , show that \(u = u' \circ j\) , where \(j: A \to A\{S^{-1}\}\) is the canonical mapping and \(u'\) is continuous; moreover \(u'\) is determined uniquely.
\item
  Let \(S_{1}\) , \(S_{2}\) be two multiplicative subsets of A and let \(S_{2}^{\prime}\) be the canonical image of \(S_{2}\) in \(A\{S^{-1}\}\) ; define a canonical isomorphism of \(A\{(S_{1}S_{2})^{-1}\}\) onto \(A\{S_{1}^{-1}\}\{S_{2}'^{-1}\}\) .
\item
  Let \(\mathfrak{a}\) be an open ideal of A and let \(\mathfrak{a}\{\mathrm{S}^{-1}\}\) be the Hausdorff completion of \(\mathrm{S}^{-1}\mathfrak{a}\) with respect to the topology induced by that on \(\mathrm{S}^{-1}\mathrm{A}\) ; show that \(\mathfrak{a}\{\mathrm{S}^{-1}\}\) is canonically identified with an open ideal of \(\mathrm{A}\{\mathrm{S}^{-1}\}\) and that the discrete ring \(\mathrm{A}\{\mathrm{S}^{-1}\} / \mathfrak{a}\{\mathrm{S}^{-1}\}\) is isomorphic to \(\mathrm{S}^{-1}(\mathrm{A} / \mathfrak{a})\) . Conversely, if \(\mathfrak{a}'\) is an open ideal of \(\mathrm{A}\{\mathrm{S}^{-1}\}\) , its inverse image \(\mathfrak{a}\) in A is an open ideal such that \(\mathfrak{a}' = \mathfrak{a}\{\mathrm{S}^{-1}\}\) . In particular, the mapping \(\mathfrak{p} \mapsto \mathfrak{p}\{\mathrm{S}^{-1}\}\) is an increasing bijection of the set of open prime ideals of A not meeting S onto the set of open prime ideals of \(\mathrm{A}\{\mathrm{S}^{-1}\}\) .
\item
  Let \(p\) be an open prime ideal of \(A\) and let \(S = A - p\) . Show that \(A\{S^{-1}\}\) is a local ring whose residue field is isomorphic to the field of fractions of \(A / p\) .
\item
  For every element \(f \in \mathbf{A}\) , let \(\mathbf{A}_{(f)}\) denote the ring \(\mathbf{A}\{\mathbf{S}_f^{-1}\}\) , where \(\mathbf{S}_f\) is the multiplicative set of the \(f^n\) ( \(n \geqslant 0\) ). If \(f\) runs through a multiplicative subset \(\mathbf{S}\) of \(\mathbf{A}\) , the \(\mathbf{A}_{(f)}\) form a direct system of rings whose direct limit is denoted by \(\mathbf{A}_{(\mathbf{S})}\) (without any topology); define a canonical homomorphism
\end{enumerate}

\[
\mathrm{A} _ {\{\mathrm{S} \}} \rightarrow \mathrm{A} \{\mathrm{S} ^ {- 1} \}.
\]

If \(S = A - p\) , where \(p\) is an open prime ideal of \(A\) , show that \(A_{(S)}\) is a local ring, the canonical homomorphism \(A_{(S)} \to A\{S^{-1}\}\) is local and the residue fields of \(A_{(S)}\) and \(A\{S^{-1}\}\) are canonically isomorphic (cf.~Chapter II, § 3, Exercise 16).

\begin{enumerate}
\def\labelenumi{(\alph{enumi})}
\setcounter{enumi}{7}
\tightlist
\item
  Suppose that the topology on A is the m-adic topology for some ideal m of A, that A is Hausdorff and complete and that \(m/m^{2}\) is a finitely generated (A/m)-module. Show that, if \(m' = m\{S^{-1}\}\) , the topology on \(A' = A\{S^{-1}\}\) is the \(m'\) -adic topology, \(m' = mA'\) and \(m'/m'^{2}\) is a finitely generated ( \(A'/m'\) )-module (use Proposition 14 of no. 10). If A is Noetherian, so is \(A'\) .
\end{enumerate}

\begin{enumerate}
\def\labelenumi{\arabic{enumi}.}
\setcounter{enumi}{27}
\tightlist
\item
  Let A be a commutative ring with a linear topology and E, F two topological A-modules which are linearly topologized. If V (resp. W) runs through the set of open submodules of E (resp. F), the submodules
\end{enumerate}

\[
\operatorname{Im} (\mathrm{V} \otimes_ {\mathbf {A}} \mathrm{F}) + \operatorname{Im} (\mathrm{E} \otimes_ {\mathbf {A}} \mathrm{W})
\]

of \(\mathbf{E} \otimes_{\mathbf{A}} \mathbf{F}\) form a fundamental system of neighbourhoods of 0 in \(\mathbf{E} \otimes_{\mathbf{A}} \mathbf{F}\) for a topology which is compatible with its module structure over the topological ring A and is called the tensor product of the given topologies on E and F. The Hausdorff completion \((\mathbf{E} \otimes_{\mathbf{A}} \mathbf{F})^{\wedge}\) of this A-module is an \(\hat{\mathbf{A}}\) -module called the completed tensor product of E and F.

\begin{enumerate}
\def\labelenumi{(\alph{enumi})}
\tightlist
\item
  Show that, if \((\mathrm{V}_{\lambda})\) (resp. \((\mathrm{W}_{\mu})\) ) is a fundamental system of neighbourhoods of 0 in E (resp. F) consisting of submodules, \((\mathbf{E} \otimes_{\mathbf{A}} \mathbf{F})^{\wedge}\) is canonically isomorphic to the inverse limit of the inverse system of A-modules
\end{enumerate}

\[
\left(\mathrm{E} / \mathrm{V} _ {\lambda}\right) \otimes_ {\mathrm{A}} \left(\mathrm{F} / \mathrm{W} _ {\mu}\right);
\]

deduce that \((\mathbf{E} \otimes_{\mathbf{A}} \mathbf{F})^{\wedge}\) is an \(\hat{\mathbf{A}}\) -module canonically isomorphic to \((\hat{\mathbf{E}} \otimes_{\hat{\mathbf{A}}} \hat{\mathbf{F}})\) ; it is also denoted by \(\hat{\mathbf{E}} \widehat{\otimes}_{\hat{\mathbf{A}}} \hat{\mathbf{F}}\) .

\begin{enumerate}
\def\labelenumi{(\alph{enumi})}
\setcounter{enumi}{1}
\item
  Let \(\mathbf{E}'\) , \(\mathbf{F}'\) be two topological A-modules which are linearly topologized and \(u: \mathbf{E} \to \mathbf{E}', v: \mathbf{F} \to \mathbf{F}'\) two continuous A-linear mappings; show that \(u \otimes v: \mathbf{E} \otimes \mathbf{F} \to \mathbf{E}' \otimes \mathbf{F}'\) is continuous with the tensor product topologies on \(\mathbf{E} \otimes \mathbf{F}\) and \(\mathbf{E}' \otimes \mathbf{F}'\) ; we denote by \(u \hat{\otimes} v\) the continuous linear mapping \((\mathbf{E} \otimes \mathbf{F})^{\wedge} \to (\mathbf{E}' \otimes \mathbf{F}')^{\wedge}\) corresponding to \(u \otimes v\) .
\item
  Let B, C be two commutative A-algebras with linear topologies such that the canonical mappings \(\mathbf{A} \to \mathbf{B}\) , \(\mathbf{A} \to \mathbf{C}\) are continuous (to abbreviate, B, C are called topological A-algebras). Show that \((\mathbf{B} \otimes_{\mathbf{A}} \mathbf{C})^{\wedge}\) has a canonical topological A-algebra structure called the completed tensor product of the algebras B and C. Define canonical continuous representations
\end{enumerate}

\[
\rho \colon \mathbf {B} \rightarrow (\mathbf {B} \otimes_ {\mathbf {A}} \mathbf {C}) ^ {\wedge}, \quad \sigma \colon \mathbf {C} \rightarrow (\mathbf {B} \otimes_ {\mathbf {A}} \mathbf {C}) ^ {\wedge}
\]

with the following property: for every complete Hausdorff commutative topological A-algebra D and every ordered pair of continuous A-homomorphisms \(u\colon B\to D\) , \(v\colon C\to D\) , there exists a unique continuous A-homomorphism \(w\colon(B\otimes_{A}C)^{\wedge}\to D\) such that \(u=w\circ\rho\) and \(v=w\circ\sigma\) .

\begin{enumerate}
\def\labelenumi{(\alph{enumi})}
\setcounter{enumi}{3}
\tightlist
\item
  Let m be an ideal of A such that the topology on A is the m-adic topology. If E and F are each given the m-adic topology, show that on \(E \otimes_{A} F\) the tensor product of the topologies on E and F is the m-adic topology. Deduce that \((\mathbf{E} \otimes_{\mathbf{A}} \mathbf{F})^{\wedge}\) is isomorphic to the inverse limits \(\varprojlim ((\mathrm{E} / \mathfrak{m}^{n+1}) \otimes_{\mathbf{A}} \mathbf{F})\) and \(\varprojlim (\mathbf{E} \otimes_{\mathbf{A}} (\mathbf{F} / \mathfrak{m}^{n+1}\mathbf{F}))\) .
\end{enumerate}

\begin{enumerate}
\def\labelenumi{\arabic{enumi}.}
\setcounter{enumi}{28}
\item
  Let A be a commutative ring, m an ideal of A such that \(\bigcap_{n\geq0}m^{n}=0\) and c is an element of A which is not a divisor of 0. Show that the transporters \(q_{n}=m^{n}:Ac\) are open under the m-adic topology and have zero intersection. If A is strictly linearly compact with the m-adic topology (Exercise 22(b)), \((q_{n})\) is a fundamental system of neighbourhoods of 0 in A for this topology.
\item
  Let K be an infinite commutative field and A the ring of formal power series \(K[[X,Y]]\) in two indeterminates, which is a complete Hausdorff Noetherian local ring. Define a multiplicative subset S of A such that \(S^{-1}A\) is not a semi-local ring (if \((\lambda_{n})\) is an infinite sequence of distinct elements of K, consider the prime ideals \(p_{n}=A(X+\lambda_{n}Y)\) of A).
\item
  Show that if A is a semi-local ring, so is the ring of formal power series A{[}{[}X{]}{]} (consider the Jacobson radical of this ring).
\end{enumerate}

\exercisesection{3}

\begin{enumerate}
\def\labelenumi{\arabic{enumi}.}
\tightlist
\item
  \begin{enumerate}
  \def\labelenumii{(\alph{enumii})}
  \tightlist
  \item
    Let K be a commutative field of characteristic p \textgreater{} 0 and B the ring of formal power series with coefficients in K in two infinite systems of indeterminates \((\mathbf{X}_{n})\) , \((\mathbf{Y}_{n})\) with only a finite number of terms of given degree ( \(\S 2\) , Exercise 12); B is a Hausdorff local ring with the m-adic topology, m being its maximal ideal. Let b be the closed ideal of B generated by the monomials \(Y_{i}X_{j}\) for \(i \geqslant 0\) , \(0 \leqslant j \leqslant i\) ; let A be the local ring B/b, which is Hausdorff with the n-adic topology, where n = m/b is its maximal ideal. Let c be the element of A equal to the class mod. b of \(\sum_{n=1}^{\infty} X_{n}^{p^{n}}\) ; show that, for \(k \geqslant 3\) , \(n^{k} \cap (Ac) \not\subset n^{2}c\) (and a fortiori the relation
  \end{enumerate}
\end{enumerate}

\[
\mathfrak {n} ^ {2} (\mathfrak {n} ^ {2 k} \cap (\mathrm{A} c)) = \mathfrak {n} ^ {2 k + 2} \cap (\mathrm{A} c)
\]

cannot hold).

\begin{enumerate}
\def\labelenumi{(\alph{enumi})}
\setcounter{enumi}{1}
\tightlist
\item
  Let \(C'\) be the commutative ring with underlying set \(A \times A\) , with multiplication \((a, x)(a', x') = (aa', ax' + a'x)\) ; let C be the subring \(A \times (Ac)\) of \(C'\) ; C and \(C'\) are local rings, whose maximal ideals we denote by r and \(r'\) , and \(C'\) is a finitely generated C-module. Show that \(r = r' \cap C\) but that the topology induced on C by the \(r'\) -adic topology is not the r-adic topology.
\end{enumerate}

\begin{enumerate}
\def\labelenumi{\arabic{enumi}.}
\setcounter{enumi}{1}
\tightlist
\item
  Let A be a Noetherian integral domain which is not a field, K its field of fractions and m an ideal \(\neq\) A of A. The m-adic topology on A is not induced by the m-adic topology on K and the latter is not Hausdorff.
\end{enumerate}

* 3. If A is the ring of a non-discrete valuation of height 1 (Chapter VI) and m is its maximal ideal, A is not Hausdorff with the m-adic topology, the closure of \{0\} being m. *

¶ 4. Let A be a semi-local ring and r its Jacobson radical. For A to be Noetherian, it is necessary and sufficient that every ideal of A be closed under the r-adic topology and that every maximal ideal of A be finitely generated. (If these conditions hold, show first that the Hausdorff completion \(\hat{A}\) of A is Noetherian, using § 2, no. 13, Corollary to Proposition 19 and no. 10, Corollary 5 to Theorem 1. Then observe that the r-adic topology is Hausdorff on A and that there exists an increasing injection of the set of ideals of A into the set of ideals of \(\hat{A}\).)

¶ 5. Let A be a commutative Noetherian ring and m its Jacobson radical.

\begin{enumerate}
\def\labelenumi{(\alph{enumi})}
\item
  For A to be strictly linearly compact with a linear topology \(\mathcal{T}\) ( \(\S\) 2, Exercise 19), it is necessary and sufficient that A be semi-local and complete with the m-adic topology; \(\mathcal{T}\) is then necessarily identical with this latter topology. (Use \(\S\) 2, Exercise 21 (d) and 22 (a) and also \(\S\) 3, no. 3, Proposition 6.)
\item
  For A to be linearly compact with a linear topology T, it is necessary and sufficient that A be semi-local and complete with the m-adic topology and that T be finer than this latter topology. (To see that it is necessary, reduce it first to the case where A is a local ring, using Exercise 21 (d) of § 2. Consider first the case where T is the discrete topology and note that A is then linearly compact with the m-adic topology. In the general case, reduce it to the case where T is a minimal topology (§ 2, Exercise 18 (c)); show that A is then strictly linearly compact and use criterion (γ) of § 2, Exercise 19 (a) and Exercise 18 (b) of § 2.)
\end{enumerate}

\begin{enumerate}
\def\labelenumi{\arabic{enumi}.}
\setcounter{enumi}{5}
\item
  Let A be a Noetherian local integral domain, m its maximal ideal and K its field of fractions; suppose that A is complete with the m-adic topology. Consider on A a linear topology T which is finer than the m-adic topology. Show that, if \(T'\) is the linear topology on K for which a fundamental system of neighbourhoods of 0 consists of the neighbourhoods of 0 in A with the topology T, \(T'\) is compatible with the field structure on K and K is complete with the topology \(T'\) .
\item
  Let A be a commutative Noetherian ring and E an A-module. Let A be given the topology \(\mathcal{T}_{m}(\mathrm{A})\) (§ 2, Exercise 24) and let \(\mathcal{T}_{m}(\mathrm{E})\) denote the linear topology on E for which a fundamental system of neighbourhoods of 0 consists of the submodules a.E, where a runs through a fundamental system of neighbourhoods of 0 for \(\mathcal{T}_{m}(\mathrm{A})\) , consisting of ideals of A. Show that if E is finitely generated, E is a Hausdorff topological A-module (with \(\mathcal{T}_{m}(\mathrm{A})\) and \(\mathcal{T}_{m}(\mathrm{E})\) ) where every submodule is closed. A fundamental system of neighbourhoods of 0 for \(\mathcal{T}_{m}(\mathrm{E})\) then consists of the submodules F of E such that E/F is a module of finite length and, for every submodule M of E, the topology \(\mathcal{T}_{m}(\mathbf{M})\) is induced by \(\mathcal{T}_{m}(\mathbf{E})\) .
\end{enumerate}

¶ 8. Let A be a Noetherian semi-local ring, n its nilradical (which is the largest nilpotent ideal of A) and m its Jacobson radical. Show that, if A (with the m-adic topology) is such that A/n is complete, then A is complete. (Reduce it to the case where n is generated by a single element c such that \(c^{2} = 0\) ; using Exercise 9 of General Topology, Chapter III, §3, reduce it to showing that n = Ac is complete and use the fact that n is an (A/n)-module.)

¶ 9. (a) Let A be a Zariski ring, m a defining ideal of the topology on A and B a ring such that \(A \subset B \subset \hat{A}\) , which is a finitely generated A-module; show that, if mB is open in B under the topology induced by the topology on \(\hat{A}\) , of necessity B = A (use Nakayama's Lemma).

* (b) Let A be a commutative Noetherian ring and m an ideal of A. Show that, if the Hausdorff completion \(\hat{\mathbf{A}}\) of A with respect to the m-adic topology is a finitely generated A-module, A is complete with the m-adic topology (reduce it to the case where A is Hausdorff; use Chapter V, § 2, no. 1, Proposition 1 and Theorem 1 to show that A is a Zariski ring with the m-adic topology and conclude with the aid of (a)). *

\begin{enumerate}
\def\labelenumi{\arabic{enumi}.}
\setcounter{enumi}{9}
\item
  Let A be a Zariski ring, m a defining ideal of the topology on A and E a finitely generated A-module. Show that, if \(\hat{E}\) is an \(\hat{A}\) -module admitting a system of generators with r elements, then the A-module E admits a system of generators with r elements (note that E/mE and \(\hat{E}/m\hat{E}\) are isomorphic and use Chapter II, § 3, Corollary 2 to Proposition 4). Is the result valid if E is not assumed to be finitely generated (cf.~Exercise 9)? Is it valid if we only assume that A is Noetherian (take A = Z)?
\item
  Let A be the ring \(\mathbf{Z}_p\) of \(p\) -adic integers ( \(p\) prime) with the \(p\) -adic topology with which it is a complete Zariski ring. Let E be the A-module \(\mathrm{A}^{(N)}\) with the \(p\) -adic topology.
\end{enumerate}

\begin{enumerate}
\def\labelenumi{(\alph{enumi})}
\item
  Show that the completion \(\hat{\mathbf{E}}\) of \(\mathbf{E}\) is identified with the submodule of \(\mathbf{A}^{\mathbf{N}}\) consisting of the sequences \((a_{n})_{n\in \mathbf{N}}\) of elements of \(\mathbf{A}\) such that \(\lim a_n = 0\) .
\item
  Let \(e_n\) be the \(n\) -th vector of the canonical basis of \(E\) . Consider in \(E\) the submodule \(F\) generated by the vectors \(e_{2n-1}\) and the submodule \(G\) generated by the vectors \(p^{2n}e_{2n} - e_{2n-1}\) ( \(n \geqslant 1\) ). Show that the topologies induced on \(F\) and \(G\) by the \(p\) -adic topology on \(E\) are the \(p\) -adic topologies on \(F\) and \(G\) but that, in \(\hat{E}\) ,
\end{enumerate}

\[
\overline {{{\mathbf {F} + \mathbf {G}}}} \neq \overline {{{\mathbf {F}}}} + \overline {{{\mathbf {G}}}}.
\]

\begin{enumerate}
\def\labelenumi{(\alph{enumi})}
\setcounter{enumi}{2}
\item
  In \(\hat{\mathbf{E}}\) , let \(a_{r} = \sum_{n=0}^{\infty} p^{n} e_{r+2}\) ( \(r \geqslant 0\) ); let \(\mathbf{H}\) be the submodule of \(\hat{\mathbf{E}}\) generated by the \(a_{r}\) ( \(r \geqslant 0\) ). Show that on H the topology induced by the \(p\) -adic topology on \(\hat{\mathbf{E}}\) is the \(p\) -adic topology on H but that \(\overline{\mathbf{E} \cap \mathbf{H}} \neq \bar{\mathbf{E}} \cap \bar{\mathbf{H}}\) .
\item
  Let L be the submodule of E generated by the \(p^{n}e_{n}\) ; show that on L the topology induced by the p-adic topology on E is distinct from the p-adic topology on L.
\end{enumerate}

\begin{enumerate}
\def\labelenumi{\arabic{enumi}.}
\setcounter{enumi}{11}
\tightlist
\item
  \begin{enumerate}
  \def\labelenumii{(\alph{enumii})}
  \tightlist
  \item
    Let A be a commutative Noetherian ring and \(m_{1}\) , \(m_{2}\) two ideals of A contained in the Jacobson radical of A; let \(A_{i}\) be the completion of A with respect to the \(m_{i}\) -adic topology. If \(m_{1} \subset m_{2}\) , show that the identity mapping on A can be extended by continuity to an injective representation of \(A_{1}\) into \(A_{2}\) (cf.~General Topology, Chapter III, § 3, no. 5, Proposition 9).
  \end{enumerate}
\end{enumerate}

\begin{enumerate}
\def\labelenumi{(\alph{enumi})}
\setcounter{enumi}{1}
\tightlist
\item
  Let \(n = A_1 m_2\) . Show that, if \(A_1\) is canonically identified with a subring of \(A_2\) , \(A_2\) is identified with the completion of \(A_1\) with respect to the \(n\) -adic topology.
\end{enumerate}

\begin{enumerate}
\def\labelenumi{\arabic{enumi}.}
\setcounter{enumi}{12}
\tightlist
\item
  Let A be a Noetherian local ring with maximal ideal m and B a ring such that \(A \subset B \subset \hat{A}\) . Suppose that B is a Noetherian local ring with maximal ideal n (Exercise 14).
\end{enumerate}

\begin{enumerate}
\def\labelenumi{(\alph{enumi})}
\item
  Show that \(\mathfrak{n} = \mathbf{B} \cap \hat{\mathfrak{m}}\) and \(\mathbf{B} = \mathbf{A} + \mathfrak{n}\) . If further \(\mathfrak{n}^2 = \mathbf{B} \cap \hat{\mathfrak{m}}^2\) , then \(\mathfrak{n} = \mathbf{B}\mathfrak{m}\) (note that then \(\mathfrak{n} = \mathbf{B}\mathfrak{m} + \mathfrak{n}^2\) ).
\item
  Let K be a commutative field, C the polynomial ring K{[}X{]}, p the prime ideal CX, A the (Noetherian) local ring \(C_{p}\) and m its maximal ideal; the completion \(\hat{A}\) is identified with the ring of formal power series K{[}{[}X{]}{]} (no. 4, Proposition 8). Let \(u = u(X)\) be an element of \(\hat{A}\) which is transcendental over the field of rational fractions K(X) (cf.~Algebra, Chapter V, §5, Exercise 13 or Functions of a Real Variable, Chapter III, §1, Exercise 14 (a)) with no constant term; let B be the subring of \(\hat{A}\) consisting of the quotients
\end{enumerate}

\[
\mathrm{P} (\mathrm{X}, u (\mathrm{X})) / \mathrm{Q} (\mathrm{X}, u (\mathrm{X})),
\]

where P and Q are polynomials in K{[}X, Y{]} such that Q(0, 0) ≠ 0. Show that B is a Noetherian local ring containing A, whose maximal ideal n is generated by X and u(X); but on B the n-adic topology is strictly coarser than the m-adic topology.

\begin{enumerate}
\def\labelenumi{(\alph{enumi})}
\setcounter{enumi}{2}
\tightlist
\item
  With the same definitions as in (b), let \(\mathbf{B}'\) be the subring of \(\mathbf{B}\) generated by \(\mathbf{A}\) and \(u(\mathbf{X})\) ; show that \(\mathbf{B}'\) is not a local ring (note that \(\mathbf{B}' \cap \mathfrak{n}\) is a maximal ideal of \(\mathbf{B}'\) but that there are elements of \(\mathbf{B}'\) not invertible in \(\mathbf{B}'\) and not belonging to \(\mathbf{B}' \cap \mathfrak{n}\) ).
\end{enumerate}

¶ 14. Let K be a commutative field, C the ring K{[}X, Y{]}, p the maximal ideal CX + CY of C, A the local ring C \(_{p}\) , m its maximal ideal and \(\hat{A}\) the completion of A, which is identified with the ring of formal power series K{[}{[}X, Y{]}{]} (no. 4, Proposition 8). Let B be the subring of \(\hat{A}\) consisting of the formal power series of the type Xf(X, Y) + (P(Y)/Q(Y)), where \(f \in \hat{A}\) and P and Q are two polynomials of K{[}Y{]} such that Q(0) ≠ 0.

\begin{enumerate}
\def\labelenumi{(\alph{enumi})}
\item
  Show that \(\mathbf{A} \subset \mathbf{B}\) , that \(\mathbf{B}\) is a local ring whose maximal ideal \(\mathfrak{n}\) is equal to \(\mathbf{B}\mathfrak{m}\) and that \(\mathfrak{n}^k = \mathbf{B} \cap \hat{\mathfrak{m}}^k\) for every integer \(k \geqslant 1\) ; \(\mathbf{B}\) is therefore everywhere dense in \(\hat{\mathbf{A}}\) and the topology induced on \(\mathbf{B}\) by the \(\hat{\mathfrak{m}}\) -adic topology on \(\hat{\mathbf{A}}\) is the \(\mathfrak{n}\) -adic topology.
\item
  Show that in B the ideal b generated by all the elements of the form \(\mathbf{X}f(\mathbf{Y})\) , where \(f(\mathbf{Y}) \in \mathbf{K}[[\mathbf{Y}]]\) , is not finitely generated (cf.~Algebra, Chapter V, § 5, Exercise 13) and therefore B is not Noetherian, although B/n and \(\mathrm{gr}(\mathbf{B}) = \mathrm{gr}(\hat{\mathbf{A}})\) are Noetherian and the ideal n is finitely generated; \(b = B \cap \hat{A}X\) and b is the closure in B of the principal ideal (not closed) BX; finally B is not a flat A-module and \(\hat{B} = \hat{A}\) is not a flat B-module (use Chapter I, § 3, no. 5, Proposition 9).
\item
  Let \(f(\mathbf{Y})\) be an invertible formal power series in \(\mathbf{K}[[\mathbf{Y}]]\) , which is not an element of \(\mathbf{K}(\mathbf{Y})\) . If \(c\) is the ideal of \(\mathbf{B}\) generated by \(\mathbf{X}\) and \(\mathbf{X}f(\mathbf{Y})\) , show that on \(c\) the \(n\) -adic topology is strictly finer than the topology induced on \(c\) by the \(n\) -adic topology on \(\mathbf{B}\) . Show that the canonical mapping \(\hat{\mathbf{B}} \otimes_{\mathbf{B}} c \to \hat{c}\) (where \(\hat{c}\) is the completion of \(c\) with respect to the \(n\) -adic topology) is not injective (consider the images of \(f(\mathbf{Y}) \otimes \mathbf{X}\) and \(1 \otimes \mathbf{X}f(\mathbf{Y})\) ; to show that these two elements are distinct in \(\hat{\mathbf{B}} \otimes_{\mathbf{B}} c\) , consider this tensor product as a quotient of the tensor product \(\hat{\mathbf{B}} \otimes_{\mathbf{K}(\mathbf{Y})} c\) ).
\item
  Let \(f_{1}(\mathbf{Y}), f_{2}(\mathbf{Y})\) be two invertible formal power series in \(\mathbf{K}[[\mathbf{Y}]]\) such that \(1, f_{1}\) and \(f_{2}\) are linearly independent over \(\mathbf{K}(\mathbf{Y})\) . Let \(c_{1}, c_{2}\) be the principal ideals of \(\mathbf{B}\) generated respectively by \(\mathbf{X}f_{1}(\mathbf{Y})\) and \(\mathbf{X}f_{2}(\mathbf{Y})\) ; the n-adic topologies on \(c_{1}\) and \(c_{2}\) are respectively identical with the topologies induced by that on \(\mathbf{B}\) and the closures of these ideals in \(\hat{\mathbf{A}} = \hat{\mathbf{B}}\) are both identical with the principal ideal \(\hat{\mathbf{A}}\mathbf{X}\) of \(\hat{\mathbf{A}}\) . Deduce that the closure of \(c_{1} \cap c_{2}\) in \(\hat{\mathbf{A}}\) is not equal to \(\bar{c}_{1} \cap \bar{c}_{2}\) and that the closure of \(c_{1}: c_{2}\) is not equal to \(\bar{c}_{1}: \bar{c}_{2}\) .
\end{enumerate}

\begin{enumerate}
\def\labelenumi{\arabic{enumi}.}
\setcounter{enumi}{14}
\tightlist
\item
  \begin{enumerate}
  \def\labelenumii{(\alph{enumii})}
  \tightlist
  \item
    Let A be a Noetherian integral domain, m an ideal of A and \(\hat{A}\) the Hausdorff completion of A with respect to the m-adic topology. Show that, if M is a torsion-free A-module, then, for every element b which is not a divisor of 0 in \(\hat{A}\) , the homothety with ratio b in the \(\hat{A}\) -module \(\hat{A} \otimes_{A} M\) is injective. (Reduce it to the case where M is finitely generated; then there exists a maximal
  \end{enumerate}
\end{enumerate}

 free system \((m_j)_{1\leqslant j\leqslant n}\) in M and \(a\in \mathbf{A}\) such that for all \(m\in \mathbf{M}\) , \(am = \sum_{j}a_{j}m_{j}\) where \(a_{j}\in \mathbf{A}\) ; every \(x\in \hat{\mathbf{A}}\otimes_{\mathbf{A}}\mathbf{M}\) therefore satisfies \(ax = \sum_{j}b_{j}\otimes m_{j}\) where \(b_{j}\in \hat{\mathbf{A}}\) ; then use the flatness of the A-module \(\hat{\mathbf{A}}\) .

\begin{enumerate}
\def\labelenumi{(\alph{enumi})}
\setcounter{enumi}{1}
\tightlist
\item
  Let K be an algebraically closed field of characteristic 0, B the polynomial ring K{[}X, Y{]} and P(X, Y) = X(X² + Y²) + (X² - Y²). Show that the ideal BP is prime in B; consider the quotient ring A = B/BP which is a Noetherian integral domain. Let m be the maximal ideal of A, the canonical image of the maximal ideal n = BX + BY of B. Show that the Hausdorff completion \(\hat{A}\) of A with respect to the m-adic topology is not an integral domain (observe that in the ring of formal power series K{[}{[}X, Y{]}{]}, P decomposes into a product of two formal power series (``double point of an irreducible cubic'')). 
\end{enumerate}

\begin{enumerate}
\def\labelenumi{\arabic{enumi}.}
\setcounter{enumi}{15}
\tightlist
\item
  Let A be a commutative Noetherian ring, m an ideal of A and E a finitely generated A-module with the m-adic topology. For every infinite set I, let \(E_{I}\) denote the set of families \((x_{i})_{i\in I}\) of elements of E such that lim \(x_{i}=0\) with respect to the filter of complements of finite subsets of I; it is a submodule of the product A-module \(E^{I}\) .
\end{enumerate}

\begin{enumerate}
\def\labelenumi{(\alph{enumi})}
\item
  Show that, if \(0 \rightarrow E' \rightarrow E \rightarrow E'' \rightarrow 0\) is an exact sequence of finitely generated A-modules, the corresponding sequence \(0 \rightarrow E_{1}' \rightarrow E_{1} \rightarrow E_{1}'' \rightarrow 0\) is exact.
\item
  Define a canonical A-homomorphism \(A_{I} \otimes_{A} E \to E_{1}\) for every A-module E and show that it is an isomorphism (first verify this when E is free and finitely generated and then use (a)).
\item
  Deduce from (b) that \(A_{I}\) is a faithfully flat A-module.
\end{enumerate}

¶ 17. (a) Let K be a commutative field, A = K{[}{[}X₁, \ldots, Xₙ{]}{]} the ring of formal power series in n indeterminates with coefficients in K and V a vector space over K. With the notation of § 2, no. 6, Example 1, let V{[}{[}X₁, \ldots, Xₙ{]}{]} denote the vector space Vⁿ with the A-module structure defined by

\[
\left(\sum_ {\alpha} c _ {\alpha} \mathbf {X} ^ {\alpha}\right) (v _ {\alpha}) = (w _ {\alpha}),
\]

where \(w_{\alpha} = \sum_{\beta + \gamma = \alpha} c_{\beta} v_{\gamma}\) . Show that, if V admits a basis with I as indexing set, the A-module \(V[[X_1, \ldots, X_n]]\) is isomorphic to \(A^I\) if I is finite and to the A-module \(A_1\) defined in Exercise 16 if I is infinite. Deduce that \(V[[X_1, \ldots, X_n]]\) is a flat A-module and is faithfully flat if V is not reduced to 0.

\begin{enumerate}
\def\labelenumi{(\alph{enumi})}
\setcounter{enumi}{1}
\item
  Let L be an extension of the field K. Deduce from (a) that the ring of formal power series \(L[[X_{1},\ldots,X_{n}]]\) is a faithfully flat module over the ring \(K[[X_{1},\ldots,X_{n}]]\) . If L has finite rank over K, \(L[[X_{1},\ldots,X_{n}]]\) is isomorphic to \(L\otimes_{K}K[[X_{1},\ldots,X_{n}]]\) .
\item
  If \(\mathbf{L}\) is an algebraic extension of \(\mathbf{K}\) , show that the ring \(\mathbf{L}[[\mathbf{X}_1, \ldots, \mathbf{X}_n]]\) is a faithfully flat module over the ring
\end{enumerate}

\[
\mathbf {L} \otimes_ {\mathbf {K}} \mathbf {K} [ [ \mathbf {X} _ {1}, \dots , \mathbf {X} _ {n} ] ]
\]

(consider L as the direct limit of its sub-extensions of finite rank over K and use Chapter I, § 2, no. 7, Proposition 9). Deduce that the ring

\[
\mathbf {B} = \mathbf {L} \otimes_ {\mathbf {K}} \mathbf {K} [ [ \mathbf {X} _ {1}, \dots , \mathbf {X} _ {n} ] ]
\]

is a Noetherian local ring whose completion is identified with \(\mathbf{L}[[\mathbf{X}_1,\ldots ,\mathbf{X}_n]]\) . In order that \(\hat{\mathbf{B}} = \mathbf{B}\) , it is necessary and sufficient that \(n = 0\) or \([\mathbf{L}:\mathbf{K}] < +\infty\) .

¶ 18. Let A be a local ring with maximal ideal m; an A-module M is called quasi-finite if M/mM is a vector space of finite rank over the residue field \(k = \mathbf{A} / \mathfrak{m}\) . In particular, if \(\mathbf{A}\) is an integral domain, the field of fractions \(\mathbf{K}\) of \(\mathbf{A}\) is a quasi-finite \(\mathbf{A}\) -module.

\begin{enumerate}
\def\labelenumi{(\alph{enumi})}
\item
  Show that if A is Noetherian and M is a quasi-finite A-module, its Hausdorff completion \(\hat{\mathbf{M}}\) with respect to the m-adic topology is a finitely generated \(\hat{\mathbf{A}}\) -module (use §2, no. 11, Corollary 2 to Proposition 14). In particular, if A is complete and M Hausdorff with the m-adic topology, M is a finitely generated A-module.
\item
  Let B be another local ring, n its maximal ideal, \(\phi: A \rightarrow B\) a local homomorphism and M a finitely generated B-module. Show that, if B is Noetherian and M is a quasi-finite A-module, then the m-adic and n-adic topologies on M are identical. (Note that M/mM is a B-module of finite length and deduce that, if b is the annihilator of the B-module M/mM, then \(V(b) = \{n\}\) in Spec(B) and therefore \(V(mB + b) = \{n\}\) . Conclude using Chapter II, §4, no. 3, Corollary 2 to Proposition 11.)
\item
  Under the hypotheses of (b), show that, if \(M \neq 0\) , B/b is a quasi-finite A-module (note that \(M \neq nM\) and M/nM is a vector space of finite rank over k; deduce that B/n is of finite rank over k).
\end{enumerate}

¶ 19. Let A be a commutative Noetherian ring, m an ideal of A and S a multiplicative subset of A. Let A be a given the m-adic topology.

\begin{enumerate}
\def\labelenumi{(\alph{enumi})}
\item
  Show that the ring \(\mathbf{A}\{\mathbf{S}^{-1}\}\) (§ 2, Exercise 27) is a flat A-module.
\item
  Let \(S'\) be another multiplicative subset of \(A\) contained in \(S\) . Show that \(A\{S^{-1}\}\) is a flat \((A\{S'^{-1}\})\) -module (use Exercise 27 (d) of § 2).
\item
  Show that \(\mathbf{A}\{\mathbf{S}^{-1}\}\) is a flat \(\mathbf{A}_{\{\mathbf{S}\}}\) -module (§ 2, Exercise 27 (g)).
\item
  Suppose that \(S = A - p\) , where \(p\) is an open prime ideal of \(A\) . Show that \(A\{S^{-1}\}\) is a faithfully flat \(A_{(S)}\) -module and deduce that the ring \(A_{(S)}\) is Noetherian (cf.~§ 2, Exercise 27 (g)).
\end{enumerate}

\begin{enumerate}
\def\labelenumi{\arabic{enumi}.}
\setcounter{enumi}{19}
\item
  Let A be a commutative ring and m an ideal of A; let A be given the m-adic topology. Let B be a commutative topological A-algebra ( \(\S 2\) , Exercise 28); suppose that B is a Zariski ring; let n be a defining ideal of B. Show that, if M is a finitely generated A-module with the m-adic topology, then on the B-module \(B \otimes_{A} M\) the tensor product of the topology on B and the m-adic topology on M is the n-adic topology; then the completion tensor product \((\mathbf{B} \otimes_{\mathbf{A}} \mathbf{M})^{\wedge}\) is isomorphic to \(\hat{B} \otimes_{A} M\) .
\item
  Let A be a commutative Noetherian ring, m an ideal of A and M and N two finitely generated A-modules.
\end{enumerate}

\begin{enumerate}
\def\labelenumi{(\alph{enumi})}
\item
  Suppose that M is Hausdorff with the m-adic topology. Show that in \(\operatorname{Hom}_{\mathbf{A}}(\mathbf{M}, \mathbf{N})\) with the m-adic topology the set of injective homomorphisms is open (use no. 1, Proposition 2 and the Artin-Rees Lemma).
\item
  Suppose that A is a complete Zariski ring and m a defining ideal of
\end{enumerate}

A. For every integer \(i\) , let \(\mathbf{A}_i = \mathbf{A} / \mathfrak{m}^{i+1}\) , \(\mathbf{M}_i = \mathbf{M} / \mathfrak{m}^{i+1}\mathbf{M}\) , \(\mathbf{N}_i = \mathbf{N} / \mathfrak{m}^{i+1}\mathbf{N}\) ; show that the topological A-module \(\operatorname{Hom}_{\mathbf{A}}(\mathbf{M}, \mathbf{N})\) is isomorphic to

\[
\varprojlim \operatorname{Hom} _ {\mathbf {A}} (\mathbf {M} _ {i}, \mathbf {N} _ {i}).
\]

Deduce that in \(\operatorname{Hom}_{\mathbf{A}}(\mathbf{M},\mathbf{N})\) the set of surjective homomorphisms is open.

¶ 22. (*) Let A be a ring, commutative or not; all the A-modules to be considered are left A-modules. Let P be a property such that: (α) if f: M → N is an injective A-module homomorphism and N has property P, then M has property P; (β) the direct sum of two A-modules with property P has property P.

\begin{enumerate}
\def\labelenumi{(\alph{enumi})}
\item
  Let M be an A-module. Show that the submodules \(M'\) of M such that \(M/M'\) has property P form a fundamental system of neighbourhoods of 0 for a linear topology \(\mathcal{T}_{P}(M)\) on M. Show that \(\mathcal{T}_{P}(A_{s})\) is compatible with the ring structure on A and that M with the topology \(\mathcal{T}_{P}(M)\) is a topological module over A with \(\mathcal{T}_{P}(A_{s})\) . Every A-module homomorphism \(f: M \to N\) is continuous with the topologies \(\mathcal{T}_{P}(M)\) and \(\mathcal{T}_{P}(N)\) .
\item
  Suppose that the following condition holds: \((\gamma)\) if \(N\) is a submodule of \(M\) with property \(P\) and, for every submodule \(L \neq 0\) of \(M\) , \(N \cap L \neq 0\) , then \(M\) has property \(P\) . Show that with these conditions, if \(F\) is a submodule of an A-module \(E\) , the topology \(\mathcal{T}_P(F)\) is induced by \(\mathcal{T}_P(E)\) . (Let \(F'\) be a submodule of \(F\) such that \(F/F'\) has property \(P\) ; consider a maximal element \(G\) among the submodules of \(E\) such that \(G \cap F = F'\) and show that \(E/G\) has property \(P\) .)
\end{enumerate}

¶ 23. (a) Let A be a commutative Noetherian ring and m an ideal of A. Let P\{M\} denote the following property: M is an A-module and every finitely generated submodule of M is annihilated by a power of m.

Show that conditions \((\alpha)\) , \((\beta)\) and \((\gamma)\) of Exercise 22 are fulfilled. (To prove \((\gamma)\) , reduce it to the case where M is finitely generated; for all \(a \in \mathfrak{m}\) , there exists by hypothesis \(k > 0\) such that \(a^k\mathrm{N} = 0\) ; use the fact that there exists \(r > 0\) such that \(\operatorname{Ker}(a_{\mathbf{M}}^r) \cap \operatorname{Im}(a_{\mathbf{M}}^r) = 0\) (Algebra, Chapter VIII, § 2, no. 2, Lemma 2).)

\begin{enumerate}
\def\labelenumi{(\alph{enumi})}
\setcounter{enumi}{1}
\item
  Show that, if M is a finitely generated A-module, the topology \(\mathcal{T}_{P}(\mathbf{M})\) is identical with the m-adic topology. Give an example of an A-module M for which \(\mathcal{T}_{P}(\mathbf{M})\) is strictly finer than the m-adic topology (cf.~Exercise 11).
\item
  Show that the conclusion of (a) does not extend to the case where A is a non-commutative left Noetherian ring and m is a two-sided ideal of A (consider the ring of lower triangular matrices of order 2 over a commutative field).
\end{enumerate}

¶ 24. A ring (commutative or not) A, not reduced to 0, is called a principal ideal ring if it has no divisors of zero and every left or right ideal of A is monogenous. Such a ring is left and right Noetherian.

\begin{enumerate}
\def\labelenumi{(\alph{enumi})}
\item
  Show that, for every element \(c \neq 0\) of \(\mathbf{A}\) , \(\mathbf{A} / \mathbf{A}c\) is an \(\mathbf{A}\) -module of finite length. (Observe that, if a decreasing sequence of ideals \(\mathbf{A}a_n\) of \(\mathbf{A}\) contains \(\mathbf{A}c\) , then \(c = b_na_n\) for all \(n\) and consider the right ideals \(b_na\) .)
\item
  Show that every submodule of a (left or right) free A-module is free (same argument as in Algebra, Chapter VII, § 3, Theorem 1).
\item
  In every A-module M, the set of elements of M which are not free is a submodule T of M called the torsion submodule of M (use Chapter II, §2, Exercise 14 (a)); M is called a torsion module if T = M; M is called torsion-free if T = 0.
\item
  Show that every finitely generated torsion-free A-module is free (use Chapter II, § 2, Exercise 23 (b)).
\item
  Show that every finitely generated A-module is the direct sum of a free module and a torsion module.
\item
  Let \(\mathfrak{a}\) be a two-sided ideal \(\neq 0\) of A. Show that there exists an element \(a \in \mathfrak{a}\) and an automorphism \(\sigma\) of A such that \(\mathfrak{a} = \mathrm{A}a = a\mathrm{A}\) and \(ax = \sigma(x)a\) for all \(x \in \mathrm{A}\) . (If \(b\) is \(a\) generator of the left ideal \(\mathfrak{a}\) , there exists an endomorphism \(\tau\) of A such that \(bx = \tau(x)b\) for all \(x \in \mathrm{A}\) ; show that, if \(a = ub\) is a generator of the right ideal \(\mathfrak{a}, u\) is invertible, using Algebra, Chapter VIII, § 2, Exercise 8 (b).)
\end{enumerate}

¶ 25. Let A be a principal ideal ring (Exercise 24), a a two-sided ideal ≠0 of A and a an element of a with the properties stated in Exercise 24 (f). Let P\{M\} denote the following property: M is a left A-module and every finitely generated submodule of M is annihilated by a power of a.

\begin{enumerate}
\def\labelenumi{(\alph{enumi})}
\item
   Show that the conditions \((\alpha)\) , \((\beta)\) and \((\gamma)\) of Exercise 22 are fulfilled (consider the homothety \(a_{\mathbf{M}}\) and argue as in Exercise 23 (a), observing that \(\operatorname{Ker}(a_{\mathbf{M}}^{r})\) and \(\operatorname{Im}(a_{\mathbf{M}}^{r})\) are submodules of \(\mathbf{M}\) ).
\item
  Show that every torsion A-module M (Exercise 24 (c)) is the direct sum of a submodule \(M_{a}\) which has property P and a submodule \(M_{a}^{\prime}\) such that the restriction to \(M_{a}^{\prime}\) of the homothety \(a_{M}\) is bijective (observe that, if N is a torsion A-module and \(a_{N}\) is injective, then \(a_{N}\) is bijective; for this, reduce it to the case where N is monogenous and use Exercise 24 (a) and also Algebra, Chapter VIII, § 2, no. 2, Lemma 2).
\item
  Let S be the set of elements \(s \in A\) whose canonical image in the ring A/a is invertible. Show that S is a multiplicative subset of A and that the following conditions, for any left ideal l of A, are equivalent:
\end{enumerate}

\((\alpha) \mathfrak{l} \cap \mathbb{S} \neq \varnothing; (\beta) \mathfrak{l} + \mathfrak{a} = \mathbb{A}; (\gamma) (\mathbb{A} / \mathfrak{l})_{\mathfrak{a}} = 0\) (in the notation of (b)); (δ) for all \(x \in \mathbb{A}\) , there exists \(s \in \mathbb{S}\) such that \(sx \in \mathfrak{l}\) .

Deduce that A has a (left or right) ring of fractions with respect to S (Chapter II, § 2, Exercise 22), which is a principal ideal ring and whose only non-zero two-sided ideals are generated by the canonical images of the ideals \(a^{n}\) ; moreover the canonical mapping of A to this ring of fractions is injective.

\begin{enumerate}
\def\labelenumi{(\alph{enumi})}
\setcounter{enumi}{3}
\tightlist
\item
  Suppose now that the two-sided ideal a is maximal. Show that, for every integer n \textgreater{} 0, the ring \(A/a^{n}\) is isomorphic to a matrix ring \(\mathbf{M}_{r}(\mathbf{B}_{n})\) over a completely primary ring (Algebra, Chapter VIII, § 6, Exercise 20). If \(\mathfrak{b}_n\) is the maximal ideal of \(\mathbf{B}_n\) , show that \(\mathfrak{b}_n^n = 0\) , that every (left or right) ideal of \(\mathbf{B}_n\) is of the form \(\mathfrak{b}_n^k\) and is monogenous. (Note on the one hand that
\end{enumerate}

\[
\mathbf {a} / \mathbf {a} ^ {n} = \mathbf {M} _ {r} (\mathfrak {b} _ {n})
\]

(Algebra, Chapter VIII, § 6, Exercise 5) and on the other that, for \(k < n\) , \(\mathfrak{b}_n^k / \mathfrak{b}_n^{k+1}\) is necessarily a simple \(B_n\) -module, without which \(\mathbf{M}_r(\mathfrak{b}_n^k)\) could not be a monogenous \((A / a^n)\) -module.)

Deduce that the completion \(\hat{\mathbf{A}}\) of \(\mathbf{A}\) with respect to the topology \(\mathcal{T}_{\mathbf{P}}(\mathbf{A}_s)\) is a matrix ring \(\mathbf{M}_r(\mathbf{B})\) over a ring \(\mathbf{B}\) with no divisors of zero and each of whose ideals (left or right) is a power of the same maximal two-sided ideal.

\begin{enumerate}
\def\labelenumi{\arabic{enumi}.}
\setcounter{enumi}{25}
\tightlist
\item
  Let B be a commutative ring and A a complete Noetherian semi-local subring of B. Let n be an ideal of B containing a power of the Jacobson radical of A and such that the n-adic topology on B is Hausdorff. Then show that the topology on A is induced by the n-adic topology on B (use Proposition 8 of §2, no. 7).
\end{enumerate}

¶ 27. Let A be a ring, B a commutative A-algebra which is a Zariski ring and N a finitely generated B-module.

\begin{enumerate}
\def\labelenumi{(\alph{enumi})}
\item
  Suppose that, for an ideal \(\mathfrak{I}\) of B contained in the Jacobson radical of B, the A-modules \(\mathbf{N} / \mathfrak{I}^{n + 1}\mathbf{N}\) are flat for \(n\geqslant 0\) . Show that N is a flat A-module. (If \(v\colon \mathbf{M}\to \mathbf{M}'\) is an injective homomorphism of finitely generated A-modules, it is necessary to prove that \(u = v\otimes 1:\mathbf{M}\otimes_{\mathbf{A}}\mathbf{N}\rightarrow \mathbf{M}'\otimes_{\mathbf{A}}\mathbf{N}\) is injective. Reduce it to proving that, if we write \(\mathbf{N}_n = \mathbf{N} / \mathfrak{I}^{n + 1}\mathbf{N}\) and \(u_{n} = v\otimes 1_{\mathbf{N}_{n}}\) , the homomorphism \(\lim u_n\) is injective, using the fact that the \(\mathfrak{I}\) -adic topologies on \(\mathbf{M}\otimes_{\mathbf{A}}\mathbf{N}\) and \(\mathbf{M}'\otimes_{\mathbf{A}}\mathbf{N}\) are Hausdorff.)
\item
  Let \(b\) be an element contained in the Jacobson radical of \(\mathbf{B}\) and such that the homothety with ratio \(b\) on \(\mathbf{N}\) is injective. Show that, if \(\mathbf{N} / b\mathbf{N}\) is a flat A-module, \(\mathbf{N}\) is a flat A-module (reduce it to the case in (a)).
\item
  Suppose further that A is a local ring with maximal ideal m and residue field \(k = A/m\) and that mB is contained in the Jacobson radical of B. Let P be a B-module which is a flat A-module and \(u: N \to P\) a homomorphism such that \(u \otimes 1_k: N \otimes_A k \to P \otimes_A k\) is injective. Then show that N is a flat A-module and that u is injective. (Reduce it to showing that, for every finitely generated A-module M, the homomorphism \(u \otimes 1_M: N \otimes_A M \to P \otimes_A M\) is injective; observe that the m-adic topology on \(N \otimes_A M\) is Hausdorff; then use §2, no. 8, Corollary 1 to Theorem 1, considering the commutative diagram
\end{enumerate}

\[
\begin{array}{c} \operatorname{gr} _ {0} (\mathbf {N}) \otimes_ {k} \operatorname{gr} (\mathbf {M}) \xrightarrow {\operatorname{gr} _ {0} (u) \otimes 1} \operatorname{gr} _ {0} (\mathbf {P}) \otimes_ {k} \operatorname{gr} (\mathbf {M}) \\ \Biggl \downarrow \\ \operatorname{gr} (\mathbf {N} \otimes_ {\mathbf {A}} \mathbf {M}) \xrightarrow [ \operatorname{gr} (u \otimes 1) ]{} \operatorname{gr} (\mathbf {P} \otimes_ {\mathbf {A}} \mathbf {M}) \end{array}
\]

and using the flatness of P, the vertical arrows being the canonical homomorphisms of § 2, Exercise 6.)

\exercisesection{4}

\begin{enumerate}
\def\labelenumi{\arabic{enumi}.}
\tightlist
\item
  Let A be a commutative ring which is Hausdorff and complete with a filtration \((a_{n})_{n\geq0}\) such that \(a_{0}=A\) . Let M, \(M'\) , N be three A-modules, each with the filtration induced by that on A and the topology defined by that filtration; we write \(\overline{M}=M/a_{1}M\) , \(\overline{M}'=M'/a_{1}M'\) , \(\overline{N}=N/a_{1}N\) . Let \(f: M\times M'\to N\) be a bilinear mapping and \(\bar{f}: \overline{M}\times\overline{M'}\to\overline{N}\) the \((A/a_{1})\) -bilinear mapping derived from f by taking quotients. Let \(y\in N\) , \(\bar{x}\in\overline{M}\) , \(\bar{x}'\in\overline{M'}\) be such that: (1) \(\bar{f}(\bar{x},\bar{x}')\) is the class \(\bar{y}\) of y in \(\overline{N}\) ; (2) every element of \(\overline{N}\) can be written as
\end{enumerate}

\[
\bar {f} (\overline {{{x}}}, \overline {{{z}}} ^ {\prime}) + \bar {f} (\overline {{{z}}}, \overline {{{x}}} ^ {\prime}),
\]

where \(\bar{z} \in \overline{\mathbf{M}}\) and \(\bar{z}' \in \overline{\mathbf{M}}'\) . Show that, if N is Hausdorff and M and M' are complete, there exists \(x \in \bar{x}\) and \(x' \in \bar{x}'\) such that \(f(x, x') = y\) (argue by induction as in the proof of Hensel's lemma). Under what conditions are \(x\) and \(x'\) determined uniquely?

\begin{enumerate}
\def\labelenumi{\arabic{enumi}.}
\setcounter{enumi}{1}
\tightlist
\item
  Let A be a local ring, m its maximal ideal, k = A/m its residue field and \(f: A \to k\) the canonical homomorphism. Let \(P \in A[X]\) be a monic polynomial of degree n.~We write \(B = A[X]/P \cdot A[X]\) and denote by x the canonical image of X in B.
\end{enumerate}

\begin{enumerate}
\def\labelenumi{(\alph{enumi})}
\item
  Let \(\mathbf{Q}, \mathbf{Q}'\) be two strongly relatively prime monic polynomials in \(\mathbf{A}[\mathbf{X}]\) such that \(\mathbf{P} = \mathbf{QQ}'\) . Show that the ring \(\mathbf{B}\) is the direct sum of the ideals \(\mathbf{B}. \mathbf{Q}(x)\) and \(\mathbf{B}. \mathbf{Q}'(x)\) .
\item
  Conversely, let \(B = b \oplus b'\) be a decomposition of B as a direct sum of two ideals. Show that there exist polynomials Q, \(Q'\) in A{[}X{]} satisfying the hypotheses of (a) and such that \(\mathfrak{b} = \mathrm{B}.Q(x)\) , \(\mathfrak{b}' = \mathrm{B}.Q'(x)\) . (Show first that b/mb and \(b'/mb'\) are generated by the images in B/mB of monic polynomials \(Q_{0} \in \bar{f}^{-1}(\mathfrak{b})\) , \(Q'_{0} \in \bar{f}^{-1}(\mathfrak{b}')\) such that
\end{enumerate}

\[
\bar {f} (\mathrm{P}) = \bar {f} (\mathrm{Q} _ {0}) \bar {f} (\mathrm{Q} _ {0} ^ {\prime}).
\]

Let \(r = \deg(Q_0)\) , \(s = \deg(Q_0')\) . Show that b is a free A-module with basis \(Q_0(x)\) , \(xQ_0(x)\) , \(\ldots\) , \(x^{s-1}Q_0(x)\) (apply Chapter II, § 3, no. 2, Proposition 5); we may then write \(x^sQ_0(x) = a_0Q_0(x) + a_1xQ_0(x) + \cdots + a_{s-1}x^{s-1}Q_0(x)\) , where \(a_i \in A\) ( \(0 \leqslant i \leqslant s-1\) ); show that, if we write

\[
\mathrm{Q} ^ {\prime} (\mathrm{X}) = \mathrm{X} ^ {s} - (a _ {0} + a _ {1} \mathrm{X} + \dots + a _ {s - 1} \mathrm{X} ^ {s - 1}),
\]

then \(\bar{f}(\mathbf{P}) = \bar{f}(\mathbf{Q}_0)\bar{f}(\mathbf{Q}')\) and \(\mathbf{Q}'\in \bar{f}^{-1}(\mathfrak{b}')\) . Similarly define \(\mathbf{Q}\) starting with \(\mathbf{Q}'\) and \(\mathbf{Q}_0\) and show that \(\mathbf{Q}\) and \(\mathbf{Q}'\) solve the problem, using Proposition 2 of no. 1.)

¶ 3. Let A be a local ring, m its maximal ideal, k = A/m its residue field and f: A → k the canonical homomorphism. Show that the two following conditions are equivalent:

\begin{enumerate}
\def\labelenumi{(\Alph{enumi})}
\setcounter{enumi}{7}
\item[(H)]
  For every monic polynomial \(\mathbf{P} \in \mathbf{A}[\mathbf{X}]\) and every decomposition of \(\bar{f}(\mathbf{P}) \in k[\mathbf{X}]\) as a product \(\bar{f}(\mathbf{P}) = \overline{\mathbf{Q}}\) . \(\overline{\mathbf{Q}}'\) of relatively prime monic polynomials, there exist two monic polynomials \(\mathbf{Q}\) , \(\mathbf{Q}'\) in \(\mathbf{A}[\mathbf{X}]\) such that \(\bar{f}(\mathbf{Q}) = \overline{\mathbf{Q}}\) , \(\bar{f}(\mathbf{Q}') = \overline{\mathbf{Q}}'\) and \(\mathbf{P} = \mathbf{QQ'}\) .
\item[(C)]
  Every commutative algebra over A which is a finitely generated A-module is a direct composition of A-algebras which are local rings.
\end{enumerate}

(To prove that (H) implies (C), argue as in Proposition 8 of no. 6, using Exercise 2 (a). To see that (C) implies (H), show first that, for every commutative A-algebra B which is a finitely generated A-module, every decomposition of B/mB as a direct sum of two ideals is necessarily of the form b/mb ⊕ b'/mb', where B = b ⊕ b' is a decomposition of B as a direct sum of two ideals; then use Exercise 2 (b).)

A local ring satisfying conditions (H) and (C) is called Henselian. Every complete Hausdorff local ring is Henselian. If A is Henselian and B is a commutative A-algebra which is a local ring and a finitely generated A-module, then B is Henselian.

\begin{enumerate}
\def\labelenumi{\arabic{enumi}.}
\setcounter{enumi}{3}
\tightlist
\item
  \begin{enumerate}
  \def\labelenumii{(\alph{enumii})}
  \tightlist
  \item
    Let \((\mathbf{A}_{\alpha}, \phi_{\alpha \beta})\) be a direct system of Henselian local rings, the homomorphisms \(\phi_{\alpha \beta}\) being local. Show that the local ring \(\mathbf{A} = \varinjlim \mathbf{A}_{\alpha}\) (Chapter II, § 3, Exercise 16) is Henselian (use criterion (H) of Exercise 3).
  \end{enumerate}
\end{enumerate}

\begin{enumerate}
\def\labelenumi{(\alph{enumi})}
\setcounter{enumi}{1}
\tightlist
\item
  Let K be a commutative field and L an algebraic extension of K. Deduce from (a) that the ring \(L \otimes_{K} K[[X_{1}, \ldots, X_{n}]]\) is Henselian. * (c) Let A be a Henselian local ring and B a commutative A-algebra which is integral over A (Chapter V) and is a local ring. Show that B is a Henselian ring (use (a)). *
\end{enumerate}

¶ 5. Let A be a Henselian local ring, B a (not necessarily commutative) A-algebra which is finitely generated A-module, b a two-sided ideal of B and let \(\overline{\mathbf{B}} = \mathbf{B} / \mathfrak{b}\) .

\begin{enumerate}
\def\labelenumi{(\alph{enumi})}
\item
  Show that every idempotent \(\varepsilon\) in \(\overline{\mathbf{B}}\) is the canonical image of an idempotent in \(\mathbf{B}\) (reduce it to the commutative case, considering the subalgebra of \(\mathbf{B}\) generated by a single element).
\item
  Let \((\varepsilon_{n})_{n\geqslant 1}\) be an infinite sequence of elements of \(\overline{\mathbf{B}}\) such that
\end{enumerate}

\[
\varepsilon_ {i} \varepsilon_ {j} = \delta_ {i j} \varepsilon_ {j}
\]

for every ordered pair of indices \((i,j)\) . Show that there exists in B an orthogonal sequence \((e_n)_{n\geq 1}\) of idempotents such that \(\varepsilon_{n}\) is the canonical image of \(e_n\) for all \(n\) . (Argue by induction as in Algebra, Chapter VIII, §6, Exercise 10; observe that, if \(e,e^{\prime}\) are two idempotents of B such that \(ee^{\prime} = 0\) , \(e^{\prime} - e^{\prime}e = e''\) is an idempotent such that \(ee'' = e''e = 0\) .)

\begin{enumerate}
\def\labelenumi{(\alph{enumi})}
\setcounter{enumi}{2}
\tightlist
\item
  Suppose now that b is the Jacobson radical of B. Let n be an integer and \((\varepsilon_{ij})\) ( \(1 \leqslant i \leqslant n, 1 \leqslant j \leqslant n\) ) a family of elements of \(\overline{B}\) such that \(\varepsilon_{ij}\varepsilon_{hk} = \delta_{jh}\varepsilon_{ik}\) and \(1 = \sum_{i=1}^{n} \varepsilon_{ii}\) . Show that there exists in B a family \((e_{ij})\) ( \(1 \leqslant i \leqslant n, 1 \leqslant j \leqslant n\) ) such that \(e_{ij}e_{hk} = \delta_{jh}e_{ik}\) and \(1 = \sum_{i=1}^{n} e_{ii}\) and that \(\varepsilon_{ij}\) is the canonical image of \(e_{ij}\) for every ordered pair of indices. (Use (b), Exercise 11 of Algebra, Chapter VIII, §6, and Exercise 9 of Algebra, Chapter VIII, §1.) Deduce that, if \(\overline{B}\) is isomorphic to a matrix ring \(\mathbf{M}_{r}(\overline{\mathbf{D}})\) , where \(\overline{D}\) is a not necessarily commutative field, then B is isomorphic to a matrix ring \(\mathbf{M}_{r}(\mathbf{D})\) , where D is an A-algebra which is a finitely generated A-module and whose Jacobson radical \(\mathfrak{d}\) is such that \(D/\mathfrak{d}\) is isomorphic to \(\overline{D}\) (cf.~Algebra, Chapter VIII, §1, Exercises 9 and 3). Show that, if further B is an Azumaya algebra over A (Chapter II, §5, Exercise 14), so is D.
\end{enumerate}

\begin{enumerate}
\def\labelenumi{\arabic{enumi}.}
\setcounter{enumi}{5}
\tightlist
\item
  Give an example of a non-commutative Artinian ring A, which is filtered by a sequence \((a_{n})\) of two-sided ideals such that \(a_{0} = A\) and \(a_{2} = 0\) and for which Proposition 8 of no. 6 does not hold (cf.~Algebra, Chapter VIII, §2, Exercise 6).
\end{enumerate}

¶ 7. (a) Let A be a commutative ring and m an ideal of A such that A is Hausdorff and complete with the m-adic topology and \(m/m^{2}\) is a finitely generated A-module. Show that the topology on \(A' = A\{X_{1}, \ldots, X_{r}\}\) is the \(m'\) -adic topology, where \(m' = mA'\) and that \(m'/m'^{2}\) is a finitely generated \((A'/m')\) -module (use Proposition 14 of §2, no. 11). In particular, if A is Noetherian, so is \(A'\) .

\begin{enumerate}
\def\labelenumi{(\alph{enumi})}
\setcounter{enumi}{1}
\tightlist
\item
  Let A be a commutative Noetherian ring and m an ideal of A such that A is Hausdorff and complete with the m-adic topology. Let \(u: A \rightarrow B\) be a continuous homomorphism from A to a Hausdorff commutative topological ring B making B into an A-algebra. Show that the following conditions are equivalent:
\end{enumerate}

(α) B is Noetherian, its topology is the mB-adic topology, B is complete and B/mB is a finitely generated algebra over A/m.

(β) B is topologically A-isomorphic to \(\lim B_{n}\) , where \((\mathrm{B}_{n})_{n\geqslant1}\) is an inverse system of discrete A-algebras such that the mappings \(\phi_{nm}:B_{m}\to B_{n}\) for \(m\geqslant n\) are surjective, the kernel of \(\phi_{nm}\) is \(m^{n+1}B_{m}\) and \(B_{1}\) is a finitely generated algebra over A/m.

\((\gamma)\) B is topologically A-isomorphic to a quotient of an algebra of the form \(\mathbf{A}\{\mathbf{X}_1,\dots ,\mathbf{X}_r\}\) by a closed ideal.

(To prove that \((\beta)\) implies \((\gamma)\) , use Proposition 14 of §2, no. 11 and Theorem 1 of §2, no. 8.)

\begin{enumerate}
\def\labelenumi{\arabic{enumi}.}
\setcounter{enumi}{7}
\tightlist
\item
  Let A be a linearly topologized commutative ring. The additive group \(A[[X_{1},\ldots,X_{p}]]\) is identified with the product group \(A^{N^{p}}\) and is given the product topology T.
\end{enumerate}

\begin{enumerate}
\def\labelenumi{(\alph{enumi})}
\item
  Show that \(\mathcal{T}\) is compatible with the ring structure on \(\mathbf{A}[[\mathbf{X}_1,\dots ,\mathbf{X}_p]]\) and is a linear topology for which the mappings \(f\mapsto \partial f / \partial \mathbf{X}_i\) are continuous.
\item
  For every element \(P = \sum_{e} \alpha_{e,P} X^{e}\) of \(A[[X_{1}, \ldots, X_{p}]]\) (notation of §1, no. 6) and every complete Hausdorff topological A-algebra B (§2, Exercise 28), an element \(\mathbf{x} = (x_{1}, \ldots, x_{p}) \in B^{p}\) is said to be substitutable in P if the family \((\alpha_{e,P}\mathbf{x}^{e})\) (where we have written \(x^{e} = x_{1}^{e_{1}} \ldots x_{p}^{e_{p}}\) for \(e = (e_{1}, \ldots, e_{p})\) ) converges to 0 in B with respect to the filter of complements of finite subsets of \(N^{p}\) . This family is then summable and its sum is denoted by \(P(\mathbf{x})\) . Show that the set of formal power series \(P \in A[[X_{1}, \ldots, X_{p}]]\) such that x is substitutable in P is a subring \(S_{x}\) of \(A[[X_{1}, \ldots, X_{p}]]\) and the mapping \(P \mapsto P(\mathbf{x})\) is a homomorphism from \(S_{x}\) to B.
\item
  Show that, if \(\mathbf{x}\) is substitutable in \(\mathbf{P}\) and \(\mathbf{y} = (y_1, \ldots, y_p)\) is an element of \(\mathbf{B}^p\) such that the \(y_i\) are topologically nilpotent for \(1 \leqslant i \leqslant p\) , then \(\mathbf{x} + \mathbf{y}\) is substitutable in \(\mathbf{P}\) .
\end{enumerate}

In particular, if there exists in B a neighbourhood of 0 consisting of topologically nilpotent elements, then, for all \(P \in A[[X_{1}, \ldots, X_{p}]]\) , the set D of \(x \in B^{p}\) which are substitutable in P is open and the mapping \(\mathbf{x} \mapsto \mathrm{P}(\mathbf{x})\) from D to B is continuous.

\begin{enumerate}
\def\labelenumi{(\alph{enumi})}
\setcounter{enumi}{3}
\tightlist
\item
  Suppose from now on that A is Hausdorff and complete and let \(A[[X_{1},\ldots,X_{p}]]\) be given the topology T. For a system of q formal power series \(P_{1},\ldots,P_{q}\) of \(A[[X_{1},\ldots,X_{p}]]\) to be substitutable in a formal power series
\end{enumerate}

\[
\mathrm{Q} \in \mathrm{A} [ [ \mathrm{X} _ {1}, \dots , \mathrm{X} _ {q} ] ],
\]

it is necessary and sufficient that the system \((\mathbf{P}_{1}(0),\ldots,\mathbf{P}_{q}(0))\) be substitutable in Q.

\begin{enumerate}
\def\labelenumi{(\alph{enumi})}
\setcounter{enumi}{4}
\item
  Suppose that \(\mathbf{x}\) is substitutable in each of the \(\mathrm{P}_k\) ( \(1 \leqslant k \leqslant q\) ) and that the system \((\mathrm{P}_1, \ldots, \mathrm{P}_q)\) is substitutable in \(\mathbb{Q}\) . Show that \(\mathbf{x}\) is substitutable in \(\mathbb{Q}(\mathrm{P}_1, \ldots, \mathrm{P}_q)\) , that the system \((\mathrm{P}_1(\mathbf{x}), \ldots, \mathrm{P}_q(\mathbf{x}))\) is substitutable in \(\mathbb{Q}\) and that \(\mathbb{Q}(\mathrm{P}_1(\mathbf{x}), \ldots, \mathrm{P}_q(\mathbf{x})) = ((\mathbb{Q}(\mathrm{P}_1, \ldots, \mathrm{P}_q))(\mathbf{x})\) if further one of the following hypotheses is assumed to be satisfied: ( \(\alpha\) ) there is in \(\mathbb{B}\) an ideal \(n\) such that the ordered pair \((\mathbb{B}, n)\) satisfies Hensel's conditions and the \(x_i\) ( \(1 \leqslant i \leqslant p\) ) belong to \(n\) ; ( \(\beta\) ) the formal power series \(\mathbb{Q}\) is restricted.
\item
  Take \(\mathbf{B} = \mathbf{A} = \mathbf{F}_2\) with the discrete topology, \(\mathbf{P} = \mathbf{X} + \mathbf{X}^2\) , \(\mathbf{Q} = \sum_{n=0}^{\infty} \mathbf{X}^{2^n}\) ; then \(\mathbf{P}\) is substitutable in \(\mathbf{Q}\) , \(x = 1\) is substitutable in \(\mathbf{P}\) and in \(\mathbf{Q}(\mathbf{P})\) and \(\mathbf{P}(x)\) is substitutable in \(\mathbf{Q}\) , but \(\mathbf{Q}(\mathbf{P}(x)) = 0\) and \((\mathbf{Q}(\mathbf{P}))(x) = 1\) .
\end{enumerate}
